\documentclass[10pt,a4paper]{amsart}
\usepackage[margin=17mm]{geometry}
\usepackage{listings}
\usepackage{hyperref}
%\usepackage[obeyspaces]{url}
\setlength{\parindent}{0pt}
\setlength{\parskip}{3pt}
\setlength{\emergencystretch}{3em}
\lstset{basicstyle=\ttfamily\footnotesize,breaklines=true,
  columns=fullflexible,keepspaces=true,showstringspaces=false,
  literate={²}{{\textsuperscript{2}}}1}
\DeclareMathOperator{\clip}{clip}
\DeclareMathOperator{\trunc}{trunc}
\DeclareMathOperator{\rnd}{round}
\DeclareMathOperator{\atan}{atan2}
\title{Formulas}
\date{12 October 2026}
\begin{document}
\maketitle

This is a compendium of some computations used in the code, presented in (hopefully) a  human-readable form. Unprefixed paths are repository paths; \textbf{B} is the BMS installation,
\textbf{I} is \path{B/Data/Add-On Israel}. Units: lb, lbf, ft$^2$ unless stated;
arguments of trigonometric functions are converted to radians. Examples are rounded for reading.

\tableofcontents

\section{Moon phase and rise/set}\label{sec:moon}
\subsection{Moon phase}\label{subsec:moon-phase}
\subsubsection*{Value}
Moon phase displayed in the Weather table, selected from the
cycle percentage $p$ for the campaign date.

\subsubsection*{Formula}

\[
 p=\trunc(100d/29.530588853)\bmod100.
\]

\subsubsection*{Variables}

\begin{itemize}
\setlength{\itemsep}{0pt}
\setlength{\parsep}{0pt}
\interlinepenalty=10000
\item $p$: integer cycle percentage.
\item $d=(\text{date}-2000\text{-}01\text{-}06).\text{days}$: whole days since
6 January 2000, where $\text{date}$ is the campaign calendar date.
\end{itemize}
\subsubsection*{Data}
\begin{itemize}
\setlength{\itemsep}{0pt}
\setlength{\parsep}{0pt}
\item $\text{date}$ is read from the year/month/day entries in the first
four little-endian unsigned 32-bit integers of \path{<save>.twx}; the first
integer is ignored. $d$ is computed from that date and the fixed epoch
6 January 2000 stored as \texttt{date0} in \path{lib/moon.py}.
\item $p$ is computed from $d$ using the fixed 29.530588853-day period in
\path{lib/moon.py:moonphase_crude}. Time of day is ignored; $p$ measures
cycle progress, not illuminated fraction.
\item The displayed phase is selected from $p$ using the thresholds in
\path{lib/moon.py:moonphase_crude}: new 0--4 or 96--99; waxing crescent
5--20; first quarter 21--30; waxing gibbous 31--45; full 46--55;
waning gibbous 56--70; third quarter 71--80; waning crescent 81--95.
\end{itemize}

\subsubsection*{Example}
\path{I/Campaign/weight_tests.twx}: 2026-10-10,
$d=9774$, $p=97\Rightarrow$ displayed phase \texttt{New Moon}.
\subsubsection*{Code}
\path{lib/moon.py}, \texttt{moonphase\_crude}:
\begin{lstlisting}[numbers=left,firstnumber=113,numberstyle=\tiny,xleftmargin=2em]
delta_days = (date - date0).days
phase = int(delta_days*100/29.530588853) % 100
\end{lstlisting}

\subsection{Moonrise and moonset}\label{subsec:moonrise-moonset}
\subsubsection*{Values}
Moonrise and moonset times for the theater centre on the
campaign date, displayed in the Weather table as \texttt{HH:MMz} (UTC).

\subsubsection*{Formulas}
The lunar-position series follows Meeus~\cite[Chapter~45]{meeus1991},
as referenced in \path{lib/moon.py:133}. Rise/set uses Julian date $JD$, lunar
longitude/latitude $(\lambda,\beta)$, obliquity $\epsilon$ and Greenwich
sidereal angle $\theta_G$. The formula is
\[
 \lambda=L'+10^{-6}\!\left[\sum_{j=1}^{n}E^{|m_j|}a_j
 \sin(d_jD+m_jM+u_jM'+f_jF)+A_\lambda\right],
\]
with the analogous sum for $\beta$ (without $L'$).

%$D,M,M',F,L',E$ and the
%additive terms are time polynomials/expressions in \path{lib/moon.py:moon_long_lat};
%the coefficient rows are \texttt{SIGMA\_LR\_TERMS} and \texttt{SIGMA\_B\_TERMS}.
%Coefficients are millionths of a degree.
Then
\begin{align*}
 \alpha&=\atan(\sin\lambda\cos\epsilon-\tan\beta\sin\epsilon,\cos\lambda),&
 \delta&=\arcsin(\sin\beta\cos\epsilon+\cos\beta\sin\epsilon\sin\lambda),\\
 h&=\arcsin[\sin\phi\sin\delta+\cos\phi\cos\delta\cos(\theta_G+\ell_E-\alpha)],&
 t_*&=t_i+120\,\frac{|h_i|}{|h_i|+|h_{i+1}|}\quad\text{seconds}.
\end{align*}

\subsubsection*{Variables}

\begin{itemize}
\setlength{\itemsep}{0pt}
\setlength{\parsep}{0pt}
\interlinepenalty=10000
\item $JD$: Julian date in days; $T=(JD-2451545)/36525$ is Julian centuries
since J2000.
\item $\lambda,\beta$: computed lunar ecliptic longitude and latitude,
respectively.
\item $L',D,F$: Moon's mean longitude, mean elongation and argument of
latitude, respectively; computed from $T$ using Meeus's fixed polynomial
coefficients in \texttt{moon\_long\_lat}.
\item $M,M'$: Sun's and Moon's mean anomalies, respectively; also computed
from $T$ using Meeus's fixed polynomial coefficients.
\item $E=1-0.002516T-0.0000074T^2$: time-dependent eccentricity correction
factor; the numerical coefficients are from Meeus.
\item $n,j$: retained term count and summation index ($j=1,\ldots,n$);
$n=5$ is the app's accuracy setting.
\item $a_j$: fixed longitude-series amplitude from Meeus, in millionths of
a degree, stored in \texttt{SIGMA\_LR\_TERMS}. The latitude series uses its
own fixed amplitudes in \texttt{SIGMA\_B\_TERMS}.
\item $d_j,m_j,u_j,f_j$: fixed integer multipliers from Meeus for
$D,M,M',F$, respectively, stored in each row of the corresponding series table.
\item $A_\lambda$: time-dependent additive longitude correction, computed
as \texttt{Sl\_additive}; its fixed amplitudes 3958, 1962 and 318 are from Meeus.
\item $\epsilon,\theta_G$: Earth's obliquity and Greenwich sidereal angle,
respectively; computed for the sample time from $T$ and $JD$.
\item $\alpha,\delta$: lunar right ascension and declination, respectively.
\item $\phi,\ell_E$: observer latitude (north positive) and longitude
(east positive), respectively; taken from the BMS theater centre.
\item $h$: geometric lunar altitude.
\item $i,t_i$: index and time (seconds) of the sample starting a two-minute
interval containing a horizon crossing: $h_i\leq0<h_{i+1}$ for moonrise,
or $h_i\geq0>h_{i+1}$ for moonset.
\item $t_*$: interpolated crossing time.
\end{itemize}
\subsubsection*{Data}
\begin{itemize}
\setlength{\itemsep}{0pt}
\setlength{\parsep}{0pt}
\item $JD,T$ are computed from the campaign date read from
\path{<save>.twx} and the two-minute sample times generated in
\path{lib/moon.py:moon_rise_set_times}; \texttt{date\_to\_JD} and
\texttt{date\_to\_T} in that file perform the conversions.
\item $\phi,\ell_E$ are read from \texttt{Center Latitude} and
\texttt{Center Longitude} in the active terrain's \path{NewTerrain/Theater.txt}
by \path{lib/theater_paths.py}. Missing coordinates use the
50 N, 50 W fallback in \path{lib/brief_render.py}.
\item $L',D,F,M,M',E,A_\lambda$ are computed from $T$ and the resulting
angles using Meeus's fixed coefficients, stored in
\path{lib/moon.py:moon_long_lat}.
\item $a_j,d_j,m_j,u_j,f_j$ are taken from Meeus's longitude-series
table, stored as \texttt{SIGMA\_LR\_TERMS} in \path{lib/moon.py}.
The $\beta$ series uses its own rows from \texttt{SIGMA\_B\_TERMS}
in the same file.
\item $n=5$ is the app's \texttt{accuracy} argument set in
\path{lib/brief_render.py}; $j$ enumerates the first $n$ table rows.
\item $\epsilon,\theta_G$ are computed from $T$ and $JD$ by
\texttt{eps} and \texttt{sidereal\_time\_Greenwich} in \path{lib/moon.py}.
\item $\lambda,\beta,\alpha,\delta,h$ are computed from these time,
table and observer inputs by \texttt{moon\_long\_lat},
\texttt{moon\_asc\_decl} and \texttt{moon\_angle\_above\_horizon} in
\path{lib/moon.py}. $h$ is geometric altitude, with no refraction or
parallax correction.
\item $i,t_i,t_*$ are obtained from the sampled altitudes and times by
\path{lib/moon.py:moon_rise_set_times}: indexes $i$ in $t_i$ is determined by the sign changes of $h$;
linear interpolation gives the rise/set time $t_*$.
\end{itemize}

\subsubsection*{Example}
\path{I/Campaign/weight_tests.twx}: 2026-10-10;
Israel \path{Theater.txt}: 32 N, 35 E. Rise 03:16:47.9, set 14:53:46.4 UTC
(displayed \texttt{Moonrise: 03:16z} / \texttt{Moonset: 14:53z}).
\subsubsection*{Code}
The excerpts follow $(\lambda,\beta)\to(\alpha,\delta)\to h$ and then find
the horizon crossings. In the final function, \texttt{tt} holds two-minute
UTC sample datetimes; \texttt{hh[i]} is the lunar altitude $h_i$ in degrees
at \texttt{tt[i]}. The lists \texttt{rise\_ends} and \texttt{set\_ends}
hold indices of crossing intervals; \texttt{rise} and \texttt{set} hold the
interpolated UTC datetimes. The loop variable \texttt{n} is a time-sample
counter, separate from the periodic-term count $n$ above.
The caller in \path{lib/brief_render.py:395} supplies \texttt{accuracy=5}
and west-positive longitude $-\ell_E$.

\path{lib/moon.py}, \texttt{moon\_long\_lat}: compute and return
$(\lambda,\beta)$.
\begin{lstlisting}[numbers=left,firstnumber=166,numberstyle=\tiny,xleftmargin=2em]
l = [(E**(abs(c[1])))*c[4]*sin((c[0]*D + c[1]*M + c[2]*Mp + c[3]*F)) for c in SIGMA_LR_TERMS[0:accuracy]]
b = [(E**(abs(c[1])))*c[4]*sin((c[0]*D + c[1]*M + c[2]*Mp + c[3]*F)) for c in SIGMA_B_TERMS[0:accuracy]]

Sl_additive = (3958*sin(A1)
               + 1962*sin(Lp - F)
               + 318*sin(A2))
Sb_additive = (-2235*sin(Lp)
               +382*sin(A3)
               +175*sin(A1-F)
               +175*sin(A1+F)
               +127*sin(Lp - Mp)
               -115*sin(Lp+Mp))

Sl = sum(l) + Sl_additive
Sb = sum(b) + Sb_additive

lambd = (degrees(Lp) + Sl/(10**6)) % 360
beta = (Sb/(10**6)) % 360

return(lambd, beta)
\end{lstlisting}

\path{lib/moon.py}, \texttt{moon\_asc\_decl}: convert these coordinates to
$(\alpha,\delta)$.
\begin{lstlisting}[numbers=left,firstnumber=188,numberstyle=\tiny,xleftmargin=2em]
def moon_asc_decl(date, accuracy = 60):
    (lambd, beta) = [radians(x) for x in moon_long_lat(date, accuracy = accuracy)]
    T = date_to_T(date)
    reps = radians(eps(T))
    alpha = atan2((sin(lambd)*cos(reps)-tan(beta)*sin(reps)),(cos(lambd)))
    delta = asin(sin((beta))*cos((reps))+cos((beta))*sin((reps))*sin((lambd)))
    return (degrees(alpha), degrees(delta))
\end{lstlisting}

\path{lib/moon.py}, \texttt{moon\_angle\_above\_horizon}: use
$(\alpha,\delta)$ and the observer coordinates to obtain $h$.
\begin{lstlisting}[numbers=left,firstnumber=197,numberstyle=\tiny,xleftmargin=2em]
def moon_angle_above_horizon(date, latitude, longitude, accuracy = 60):
    rlatitude = radians(latitude)
    (alpha, delta) = moon_asc_decl(date, accuracy = accuracy)
    rdelta = radians(delta)
    H = sidereal_time_Greenwich(date) - longitude - alpha
    rH = radians(H)
    return degrees(asin(sin(rlatitude)*sin(rdelta) + cos(rlatitude)*cos(rdelta)*cos(rH)))
\end{lstlisting}

\path{lib/moon.py}, \texttt{moon\_rise\_set\_times}: build \texttt{tt} and
\texttt{hh}, select sign changes, and interpolate \texttt{rise} and \texttt{set}.
\begin{lstlisting}[numbers=left,firstnumber=206,numberstyle=\tiny,xleftmargin=2em]
def moon_rise_set_times(day, latitude, longitude, accuracy = 60):
    # compute angle every two minutes
    date = dt.datetime.combine(day, dt.time.min)
    tt = [date + n*dt.timedelta(minutes = 2) for n in range(24*30)]
    hh = [moon_angle_above_horizon(t, latitude = latitude, longitude = longitude, accuracy = accuracy) for t in tt]
    rise_ends = [i for i in range(len(tt)-1) if hh[i] <= 0 and hh[i+1] > 0]
    set_ends = [i for i in range(len(tt)-1) if hh[i] >= 0 and hh[i+1] < 0]
    rise = [tt[i] + (tt[i+1] - tt[i])*abs(hh[i])/(abs(hh[i]) +  abs(hh[i+1])) for i in rise_ends]
    set = [tt[i] + (tt[i+1] - tt[i])*abs(hh[i])/(abs(hh[i]) +  abs(hh[i+1])) for i in set_ends]
    return (rise, set)
\end{lstlisting}


\section{Takeoff fuel, gross weight and displayed DRAG}\label{sec:loadout}

\subsection{Takeoff fuel}\label{subsec:takeoff-fuel}
\subsubsection*{Value}
Takeoff fuel for the selected pilot's saved aircraft, displayed as
\texttt{T/O Fuel} in the Admin section, rounded to whole lb. Unavailable values show a dash.

\subsubsection*{Formula}
\[
 F_{TO}=F_{save}+\sum_s n_s C_s+bC_{CFT}.
\]

\subsubsection*{Variables}
\begin{itemize}
\setlength{\itemsep}{0pt}
\setlength{\parsep}{0pt}
\interlinepenalty=10000
\item $F_{TO},F_{save}$: computed takeoff fuel and the saved member's signed
fuel offset, respectively, in lb.
\item $s,n_s$: occupied station index and its saved store count. The
zero-weight internal gun/ammunition entry is omitted.
\item $C_s$: fuel capacity per store, in lb; WCD \texttt{Strength}
when \texttt{Flags \& 0x02} is nonzero, otherwise zero.
\item $b$: 1 if the aircraft supports CFTs and the saved member has them
fitted, otherwise 0.
\item $C_{CFT},C_{int}$: CFT and internal fuel capacities, in lb.
\item $C_{fit}=C_{int}+\sum_s n_s C_s+bC_{CFT}$: fitted capacity, in lb;
the app requires $0\leq F_{TO}\leq C_{fit}$.
\end{itemize}

\subsubsection*{Data}
\begin{itemize}
\setlength{\itemsep}{0pt}
\setlength{\parsep}{0pt}
\item $F_{save}$ is decoded from the member's \texttt{fuel\_initial\_raw}
entry in the embedded \texttt{.uni} of \path{<save>.cam}; the array contains
four signed little-endian 32-bit offsets. This offset can be negative and
is not a literal internal-fuel quantity.
\item $s,n_s$ come from the loadout's \texttt{weapon\_ids} and
\texttt{weapon\_counts} in that \texttt{.uni} file. A single loadout is
shared; with one per active member, each uses its own. Other loadout
counts are rejected.
\item $C_s$ is taken from \texttt{Strength} when \texttt{Flags \& 0x02}
is nonzero in the theater's \path{Terrdata/objects/Falcon4_WCD.xml},
otherwise it is zero. The saved store ID selects the WCD \texttt{Num} row.
\item $b$ is computed from the member's byte flag in
\texttt{loaded\_cft\_raw} in the CAM's embedded \texttt{.uni}, together
with \path{aero_new2.has_cft} from the aircraft's \path{Sim/Acdata/*.txtpb}.
The saved flag is used only when the aircraft supports CFTs.
\item $C_{CFT},C_{int}$ are read from \path{auxaero.fuel.cft_fuel} and
\path{airframe.internal_fuel_lbs}, respectively, in the aircraft's
\path{Sim/Acdata/*.txtpb}.
\item $F_{TO},C_{fit}$ are computed from the preceding saved values and
capacities in \path{lib/cam/loadout_data.py:build_aircraft_members},
after \texttt{\_loadout\_data} in that file sums $n_sC_s$.
\end{itemize}

\subsubsection*{Example}
Falcon2-1, flight 2887/member 0 in \path{I/Campaign/weight_tests.cam},
uses \path{I/Sim/Acdata/F-16I-52+ IAF.txtpb}.
The save gives $F_{save}=5920$, $b=1$ and one TK300 tank
(WCD \texttt{Num=782}, \texttt{Flags=2}, \texttt{Strength=2040}).
The aircraft file gives $C_{CFT}=3060$ and $C_{int}=5920$:
\[
 F_{TO}=5920+1(2040)+1(3060)=11020\ \text{lb}.
\]
This equals $C_{fit}$; Admin displays \texttt{11,020}.

\subsubsection*{Code}
Here \texttt{flight} is the parsed CAM flight, \texttt{station} is the
current saved hardpoint, and \texttt{weapon} is its WCD row.
In \texttt{\_loadout\_data}, \texttt{external\_fuel} starts at zero and
accumulates $n_sC_s$ for each retained station. Its returned total is passed
to \texttt{build\_aircraft\_members}. There, \texttt{sim} is the parsed
aircraft data, \texttt{index} selects the saved member and
\texttt{loaded\_cft} is $b$. The helper \texttt{\_weight} checks that its
argument is a finite, nonnegative number.

\path{lib/cam/loadout_data.py}, \texttt{\_loadout\_data}:
\begin{lstlisting}[numbers=left,firstnumber=67,numberstyle=\tiny,xleftmargin=2em]
external_fuel += station.weapon_count * _weight(weapon.fuel_capacity_lb)
\end{lstlisting}

\path{lib/cam/loadout_data.py}, \texttt{build\_aircraft\_members}:
\begin{lstlisting}[numbers=left,firstnumber=152,numberstyle=\tiny,xleftmargin=2em]
cft_fuel = _weight(sim.cft_fuel_lb) if loaded_cft else 0
# BMS saves a signed offset, not a literal internal-fuel quantity.
# At empty fuel this offset cancels fitted tank and CFT capacities.
fuel = flight.initial_fuel_lb[index] + _weight(external_fuel) + cft_fuel
capacity = _weight(sim.airframe.internal_fuel_lb) + _weight(external_fuel) + cft_fuel
if not 0 <= fuel <= capacity:
    raise ValueError("Planned fuel is outside fitted capacity")
member["takeoff_fuel_lb"] = fuel
\end{lstlisting}

\subsection{Takeoff gross weight}\label{subsec:gross-weight}
\subsubsection*{Value}
Takeoff gross weight for the selected pilot's saved aircraft, displayed as
\texttt{Gross Wt} in Admin, rounded to whole lb. Unavailable values show a dash.

\subsubsection*{Formula}
\[
 W_{TO}=W_{empty}+F_{TO}+bW_{CFT}
       +\sum_s(n_sw_s+w_{rack,s}+w_{pylon,s}).
\]

\subsubsection*{Variables}
\begin{itemize}
\setlength{\itemsep}{0pt}
\setlength{\parsep}{0pt}
\interlinepenalty=10000
\item $W_{TO}$: computed takeoff gross weight, in lb.
\item $W_{empty},W_{CFT}$: clean-aircraft and dry CFT weights,
respectively, in lb.
\item $F_{TO},b$: takeoff fuel and fitted-CFT flag from section~\ref{subsec:takeoff-fuel}.
\item $s,n_s$: retained occupied station index and its saved store count.
\item $w_s,w_{rack,s},w_{pylon,s}$: dry weight per store, rack
and pylon on station $s$, respectively, in lb. Hardware counts once per
occupied station; an explicitly absent rack or pylon contributes zero.
\end{itemize}

\subsubsection*{Data}
\begin{itemize}
\setlength{\itemsep}{0pt}
\setlength{\parsep}{0pt}
\item $W_{empty},W_{CFT}$ are read from \path{airframe.empty_weight_lbs}
and \path{auxaero.fuel.cft_empty_weight}, respectively, in the aircraft's
\path{Sim/Acdata/*.txtpb}.
\item $F_{TO},b,s,n_s$ reuse the fuel calculation and saved member/loadout
data from \path{<save>.cam} described in section~\ref{subsec:takeoff-fuel}.
\item $w_s$ is read from \texttt{Weight} in
\path{Terrdata/objects/Falcon4_WCD.xml}, selecting \texttt{Num} by the
saved store ID. It excludes fuel contents. Internal stores add weight;
the zero-weight gun/ammunition entry at station 0 is already accounted
for in clean weight and is skipped.
\item $w_{rack,s},w_{pylon,s}$ are read from WCD \texttt{Weight} after
resolving hardware. The aircraft \texttt{.txtpb}'s
\path{auxaero.hardpoints.hardpoint_grp} selects a group in
\path{Terrdata/objects/BmsRack.dat}; the first compatible rack with
sufficient \texttt{RackStations} capacity is used. Its \texttt{RackCT}
and \texttt{PylonCT} select hardware by WCD \texttt{CtIdx}.
Compatibility uses \path{Falcon4_CT.xml} \texttt{MoverDefinitionData}
to link WCD \texttt{CtIdx} to \path{Falcon4_SWD.XML}
\texttt{Num}/\texttt{WpnClass}; these XML files share the objects directory.
Unresolved hardware leaves gross weight unavailable.
\item $W_{TO}$ is computed by summing the preceding fuel, airframe,
CFT, store and hardware contributions in
\path{lib/cam/loadout_data.py:build_aircraft_members}; the store/hardware
sum is computed by \texttt{\_loadout\_data} in that file.
\end{itemize}

\subsubsection*{Example}
For the same Falcon2-1 loadout, WCD store rows 158, 256, 782, 179 and 181
give four A120Bs, four CB52Bs, one TK300, one AAQ13 and one AAQ28:
\[
 4(338)+4(785)+392+429+441=5754\ \text{lb}.
\]
Selected hardware rows 661, 585, 780, 779 and 622 contribute
\[
 2(95)+2(369)+2(320)+172+42=1782\ \text{lb}.
\]
With $W_{empty}=22000$, $W_{CFT}=1700$, $b=1$ and $F_{TO}=11020$:
\[
 W_{TO}=22000+11020+1700+(5754+1782)=42256\ \text{lb}.
\]
Admin displays \texttt{42,256}.

\subsubsection*{Code}
In \texttt{\_loadout\_data}, \texttt{dry\_weight} starts at zero.
The resolver supplies the station's rack and pylon; the returned sum is
passed to \texttt{build\_aircraft\_members}. There, \texttt{clean} is
$W_{empty}$, \texttt{cft\_dry} is $bW_{CFT}$ and \texttt{fuel} is $F_{TO}$
from section~\ref{subsec:takeoff-fuel}. The remaining names and validation helper are as above.

\path{lib/cam/loadout_data.py}, \texttt{\_loadout\_data}:
\begin{lstlisting}[numbers=left,firstnumber=74,numberstyle=\tiny,xleftmargin=2em]
hardware = resolve_loadout_hardware(
    flight.support, aircraft.vehicle.number, station.hardpoint_index,
    station.weapon_id, station.weapon_count,
)
dry_weight += (station.weapon_count * _weight(weapon.weight_lb)
               + _weight(hardware.rack.weight_lb) + _weight(hardware.pylon.weight_lb))
\end{lstlisting}

\path{lib/cam/loadout_data.py}, \texttt{build\_aircraft\_members}:
\begin{lstlisting}[numbers=left,firstnumber=164,numberstyle=\tiny,xleftmargin=2em]
cft_dry = _weight(sim.fuel_system.get("cft_empty_weight")) if loaded_cft else 0
clean = _weight(sim.airframe.empty_weight_lb)
if clean == 0:
    raise ValueError("Missing clean weight")
member["gross_weight_lb"] = clean + _weight(dry_weight) + cft_dry + fuel
\end{lstlisting}

\subsection{Loadout drag factor}\label{subsec:drag}
\subsubsection*{Value}
The selected pilot's saved loadout drag factor, displayed as \texttt{DRAG}
in Admin as a whole, dimensionless number. Unavailable values show a dash.

\subsubsection*{Formula}
\[
 DI=1+bDI_{CFT}
       +\sum_{s\;\mathrm{external}}(n_sDI_s+DI_{rack,s}+DI_{pylon,s}).
\]

\subsubsection*{Variables}
\begin{itemize}
\setlength{\itemsep}{0pt}
\setlength{\parsep}{0pt}
\interlinepenalty=10000
\item $DI$: computed displayed drag factor; 1 is the fixed clean-aircraft
display baseline used by the app and checked against BMS.
\item $b,DI_{CFT}$: fitted-CFT flag from section~\ref{subsec:takeoff-fuel} and the aircraft's
CFT drag increment, respectively.
\item $s,n_s$: occupied external station index and its saved store count.
A station is external when its \texttt{is\_external\_hardpoint}
flag is 1.
\item $DI_s,DI_{rack,s},DI_{pylon,s}$: fixed WCD drag indices per store,
rack and pylon, respectively. Hardware counts once per external station.
\end{itemize}

\subsubsection*{Data}
\begin{itemize}
\setlength{\itemsep}{0pt}
\setlength{\parsep}{0pt}
\item $DI_s,DI_{rack,s},DI_{pylon,s}$ are read from \texttt{Drag} in
\path{Terrdata/objects/Falcon4_WCD.xml}; store IDs and hardware selection
use the same files and rules as section~\ref{subsec:gross-weight}.
\item $s,n_s$ come from the saved loadout in \path{<save>.cam}, restricted
to stations whose \path{auxaero.hardpoints.is_external_hardpoint.values}
entry is 1 in the aircraft's \path{Sim/Acdata/*.txtpb}.
Internal stores and their hardware contribute no drag to the sum.
\item $b$ is obtained from the saved CFT flag and aircraft support flag
as in section~\ref{subsec:takeoff-fuel}; $DI_{CFT}$ is read from \path{auxaero.fuel.cft_drag}
in the aircraft's \path{Sim/Acdata/*.txtpb}. This increment is added
only when CFTs are fitted.
\item The baseline 1 in $DI$ is a fixed display convention in
\path{lib/cam/loadout_data.py:build_aircraft_members}, checked against
BMS's clean-loadout display. It is not an aerodynamic drag coefficient
or an aircraft-file field.
\item $DI$ is computed from the external store/hardware sum, baseline
and fitted CFT increment in \path{lib/cam/loadout_data.py}.
Missing drag data leaves DRAG unavailable independently of otherwise
known fuel and weight.
\end{itemize}

\subsubsection*{Example}
All retained Falcon2-1 stations in this save are external. The same store
and hardware WCD rows listed in section~\ref{subsec:gross-weight} give, respectively,
\[
 4(4)+4(22)+17+22+22=165,\qquad
 2(2)+2(24)+2(14)+16+3=99.
\]
The aircraft file gives $DI_{CFT}=20$, hence
\[
 DI=1+1(20)+165+99=285.
\]
Admin displays \texttt{285}, matching the BMS loadout-screen reading.
Removing CFTs gives 265; a clean loadout without CFTs gives 1.

\subsubsection*{Code}
In \texttt{\_loadout\_data}, \texttt{drag\_index} starts at zero and accumulates
the external store/hardware sum; \texttt{external} contains the
station flags. The \texttt{hardware} objects were resolved in section~\ref{subsec:gross-weight}.
The helper \texttt{\_drag\_index} validates a nonnegative integer.
In \texttt{build\_aircraft\_members}, the returned raw sum is augmented
with the clean baseline and fitted CFT increment before assignment to
\texttt{member["drag\_index"]}.

\path{lib/cam/loadout_data.py}, \texttt{\_loadout\_data}:
\begin{lstlisting}[numbers=left,firstnumber=85,numberstyle=\tiny,xleftmargin=2em]
external = aircraft.simulation_data.external_hardpoints
if not 0 <= station.hardpoint_index < len(external):
    raise ValueError("Missing external hardpoint flag")
flag = external[station.hardpoint_index]
if flag not in (0, 1):
    raise ValueError("Invalid external hardpoint flag")
if flag:
    # Raw loadout DI: each store, plus one rack/pylon per station.
    # Internal stores are excluded. The displayed baseline and
    # fitted CFT increment are added per member below.
    drag_index += (station.weapon_count * _drag_index(weapon.drag_index)
                   + _drag_index(hardware.rack.drag_index)
                   + _drag_index(hardware.pylon.drag_index))
\end{lstlisting}

\path{lib/cam/loadout_data.py}, \texttt{build\_aircraft\_members}:
\begin{lstlisting}[numbers=left,firstnumber=148,numberstyle=\tiny,xleftmargin=2em]
cft_drag = _drag_index(sim.fuel_system.get("cft_drag")) if loaded_cft else 0
member["drag_index"] = 1 + _drag_index(drag_index) + cft_drag
\end{lstlisting}

\section{Atmosphere, aerodynamic drag and fuel flow}\label{sec:flow}
This is a technical section required for further calculations. None of the values are displayed in the briefing directly.

\subsection{Atmosphere}\label{subsec:atmosphere}
\subsubsection*{Values}
Mach number, true airspeed and dynamic pressure at a given MSL altitude
and calibrated airspeed. They supply the route-performance model below;
the liftoff model uses a fixed sea-level reference.

\subsubsection*{Formulas}
Temperature and pressure follow the first two layers of the
\emph{U.S. Standard Atmosphere, 1976} \cite[Part~1]{ussa1976}.
The app uses MSL altitude directly as layer height, without converting
to geopotential height. The subsonic airspeed and pressure relations
follow \cite[eqs.~(1), (2), (5)--(7) and (12)]{aiken1946}:
\begin{align*}
 M&=\sqrt{\frac{2}{\kappa-1}
       \left[\left(1+\frac{q_c}{p}\right)^{(\kappa-1)/\kappa}-1\right]},\\
 V_{TAS}&=M\sqrt{\kappa RT}\,\frac{3600}{1852},&
 q&=\frac{\kappa}{2}pM^2\,\frac{0.09290304}{4.4482216152605}.
\end{align*}

\subsubsection*{Variables}
\begin{itemize}
\setlength{\itemsep}{0pt}
\setlength{\parsep}{0pt}
\interlinepenalty=10000
\item $M,V_{TAS},q$: computed Mach number (dimensionless), true airspeed
in knots and dynamic pressure in lbf/ft$^2$, respectively.
\item $h,V_{CAS}$: input MSL altitude in feet and calibrated airspeed in
knots; $z=0.3048h$ and $v_c=V_{CAS}(1852/3600)$ are their metre and
metre-per-second equivalents.
\item $R=287.05287$ J/(kg K), $g=9.80665$ m/s$^2$, $\kappa=1.4$:
fixed dry-air gas constant, gravitational acceleration and specific-heat
ratio used by the code.
\item $T_0=288.15$ K, $p_0=101325$ Pa: fixed ISA sea-level temperature
and pressure.
\item $L_T=0.0065$ K/m, $z_t=11000$ m: fixed temperature lapse rate
below the tropopause and its height; the layer above is isothermal.
\item $T=T_0-L_T\min(z,z_t)$: computed air temperature in K.
\item $p$: computed static pressure in Pa,
\[
 p=p_0(T/T_0)^{g/(RL_T)}
       \exp\!\left[-\frac{g\max(z-z_t,0)}{RT}\right].
\]
\item $q_c$: impact pressure in Pa corresponding to the input CAS,
\[
 q_c=p_0\left[\left(1+\frac{\kappa-1}{2}
              \frac{v_c^2}{\kappa RT_0}\right)^{\kappa/(\kappa-1)}-1\right].
\]
\item $0.3048$, $1852$, $3600$: metres per foot, metres per nautical
mile and seconds per hour. The pressure conversion uses
$0.09290304$ m$^2$/ft$^2$ and $4.4482216152605$ N/lbf.
\end{itemize}

\subsubsection*{Data}
\begin{itemize}
\setlength{\itemsep}{0pt}
\setlength{\parsep}{0pt}
\item $h,V_{CAS}$ are obtained from the altitude and CAS columns in
\path{B/User/Briefings/briefing.txt}. In
\path{lib/fuel_brief.py:_altitude}, an altitude suffix \texttt{M} means
thousands of feet; AGL suffix \texttt{A} is rejected. Airport altitude
instead comes from the negative of the third coordinate in
\texttt{[STPT]} entries \texttt{target\_k} of the loaded DTC
\path{<callsign/save>.ini}. Leg midpoint altitudes and speed selection
are described in section~\ref{sec:mfob}.
\item $h=0$, $V_{CAS}=100$ are fixed inputs in \path{lib/liftoff.py}
for the liftoff reference; they are independent of the departure field.
\item $T,p,q_c,M,V_{TAS},q$ are computed from these inputs and the fixed
constants above by \path{lib/route_fuel.py:air_state}.
No saved weather is used. The helper requires $-500\leq z\leq20000$ m,
positive CAS, and both local Mach and the sea-level Mach corresponding
to CAS below 1.
\end{itemize}

\subsubsection*{Example}
The 2026-10-12 11:57:32 export of
\path{B/User/Briefings/briefing.txt}, Weasel7 STPT 3, gives altitude
\texttt{5.9M} and CAS \texttt{315}. Thus $h=5900$ ft,
$z=1798.32$ m and $v_c=162.05$ m/s. The formulas give
\[
 T=288.15-0.0065(1798.32)=276.46092\ \text{K},\qquad
 p=81506.13\ \text{Pa},\quad q_c=17017.04\ \text{Pa},
\]
\[
 M=0.527587,\qquad V_{TAS}=341.836\ \text{kt},\qquad
 q=331.681\ \text{lbf/ft}^2.
\]

\subsubsection*{Code}
The arguments \texttt{altitude\_ft} and \texttt{cas\_kt} are $h,V_{CAS}$;
after conversion, \texttt{altitude} and \texttt{cas} are $z,v_c$.
The numeric factors $0.2,3.5,5,2/7,0.7$ are the corresponding expressions
in $\kappa=1.4$. \texttt{\_number} validates finite numeric inputs.

\path{lib/route_fuel.py}, \texttt{air\_state}:
\begin{lstlisting}[numbers=left,firstnumber=17,numberstyle=\tiny,xleftmargin=2em]
def air_state(altitude_ft: float, cas_kt: float) -> tuple[float, float, float]:
    """Return Mach, TAS in knots and dynamic pressure in lbf/ft² (subsonic)."""
    altitude = _number(altitude_ft, "MSL altitude") * .3048
    cas = _number(cas_kt, "CAS", positive=True) * 1852 / 3600
    if not -500 <= altitude <= 20000:
        raise FuelModelError("MSL altitude range")
    temperature = 288.15 - .0065 * min(altitude, 11000)
    pressure = 101325 * (temperature / 288.15) ** (9.80665 / (287.05287 * .0065))
    if altitude > 11000:
        pressure *= math.exp(-9.80665 * (altitude - 11000) / (287.05287 * temperature))
    sea_mach = cas / math.sqrt(1.4 * 287.05287 * 288.15)
    impact_pressure = 101325 * ((1 + .2 * sea_mach ** 2) ** 3.5 - 1)
    mach = math.sqrt(5 * ((impact_pressure / pressure + 1) ** (2 / 7) - 1))
    if sea_mach >= 1 or mach >= 1:
        raise FuelModelError("Subsonic CAS range")
    tas = mach * math.sqrt(1.4 * 287.05287 * temperature) * 3600 / 1852
    q = .7 * pressure * mach ** 2 * .09290304 / 4.4482216152605
    return mach, tas, q
\end{lstlisting}

\subsection{Aircraft table interpolation}\label{subsec:table-interpolation}
\subsubsection*{Value}
An aircraft-table value at a requested pair of coordinates, after applying
the table's multiplier. The same helper supplies lift, drag, correction,
thrust and fuel-flow values to the performance models.

\subsubsection*{Formula}
\[
 A(x,y)=(1-u)(1-v)A_{00}+u(1-v)A_{10}
                    +(1-u)vA_{01}+uvA_{11}.
\]

\subsubsection*{Variables}
\begin{itemize}
\setlength{\itemsep}{0pt}
\setlength{\parsep}{0pt}
\interlinepenalty=10000
\item $A(x,y)$: interpolated value, in the table's units; for base lift,
$A=C_L$ is the dimensionless lift coefficient.
\item $x,y$: query coordinates, such as Mach $M$ and angle of attack
$\alpha$ in degrees for lift, or altitude in feet and Mach for engine tables.
\item $x_0,x_1,y_0,y_1$: adjacent axis breakpoints enclosing the query;
queries outside the available axes are rejected.
\item $u=(x-x_0)/(x_1-x_0)$, $v=(y-y_0)/(y_1-y_0)$: interpolation
fractions within the two intervals.
\item $A_{ij}=\mu a_{ij}$: scaled corner values. Here $a_{ij}$ are raw
table entries, $\mu$ is the multiplier (taken from BMS data; 1 if absent), and
$i,j\in\{0,1\}$ select the lower/upper corner on each axis.
\end{itemize}

\subsubsection*{Data}
\begin{itemize}
\setlength{\itemsep}{0pt}
\setlength{\parsep}{0pt}
\item $a_{ij},\mu$ are taken from a table's \path{data.values} rows and
\texttt{multiplier} in the active aircraft's \path{Sim/Acdata/*.txtpb}.
\path{lib/cam/opencam/aircraft_data.py:_table} parses them;
\path{lib/fuel.py:_Table.read} applies the multiplier once.
\item $x_0,x_1,y_0,y_1$ are selected from
\path{x_axis.breakpoints}/\path{y_axis.breakpoints} in that table.
If an axis is absent, the caller supplies fallback breakpoints from
the same aircraft file: for example, \path{aero_new2.mach_bp.breakpoints}
and \path{aero_new2.alpha_bp.breakpoints} for base lift/drag.
The selections for other tables are listed below and in
section~\ref{sec:takeoff}.
\item $x,y$ are query arguments computed from the flight state in
\path{lib/route_fuel.py} or \path{lib/liftoff.py}.
\path{lib/fuel.py:_bracket} computes the enclosing intervals and
$u,v$; \texttt{\_Table.at} in that file computes $A(x,y)$.
\end{itemize}

\subsubsection*{Example}
In \path{I/Sim/Acdata/F-16C-40 IAF.txtpb},
\path{aero_new2.clift_table} has multiplier 1. Between Mach
$0.5,0.6$ and angles $5^\circ,6^\circ$, its corner values are
\[
 (A_{00},A_{10},A_{01},A_{11})=(0.364,0.379,0.441,0.453).
\]
For a chosen query $M=0.55$, $\alpha=5.5^\circ$, $u=v=0.5$:
\[
 C_L(0.55,5.5)=\tfrac14(0.364+0.379+0.441+0.453)=0.40925.
\]

\subsubsection*{Code}
\texttt{\_bracket} returns an axis interval index and its interpolation
fraction, using the final interval at the upper endpoint.
In \texttt{read}, \texttt{table} holds the parsed raw rows and multiplier;
\texttt{x,y} are fallback axes. In \texttt{at}, they are query coordinates,
and \texttt{fx,fy} are $u,v$; \texttt{a,b} interpolate along each row.
The helpers \texttt{\_axis} and \texttt{\_number} validate the inputs.

\path{lib/fuel.py}, \texttt{\_bracket}:
\begin{lstlisting}[numbers=left,firstnumber=51,numberstyle=\tiny,xleftmargin=2em]
def _bracket(axis: tuple[float, ...], value: float) -> tuple[int, float]:
    if not axis[0] <= value <= axis[-1]:
        raise FuelModelError("Aircraft table range")
    index = min(max(bisect_right(axis, value) - 1, 0), len(axis) - 2)
    return index, (value - axis[index]) / (axis[index + 1] - axis[index])
\end{lstlisting}

\path{lib/fuel.py}, \texttt{\_Table.read/at}:
\begin{lstlisting}[numbers=left,firstnumber=65,numberstyle=\tiny,xleftmargin=2em]
def read(cls, table: AircraftTable | None, x=None, y=None) -> "_Table":
    if table is None:
        raise FuelModelError("Required aircraft performance table is missing.")
    xs = _axis(table.x_breakpoints if table.x_breakpoints is not None else x, "x")
    ys = _axis(table.y_breakpoints if table.y_breakpoints is not None else y, "y")
    if len(table.rows) != len(xs) or any(len(row) != len(ys) for row in table.rows):
        raise FuelModelError("Aircraft table dimensions do not match its breakpoints.")
\end{lstlisting}
\begin{lstlisting}[numbers=left,firstnumber=73,numberstyle=\tiny,xleftmargin=2em]
    scale = 1.0 if table.multiplier is None else _number(table.multiplier, "table multiplier", positive=True)
    rows = tuple(tuple(_number(v, "table value") * scale for v in row) for row in table.rows)
    return cls(xs, ys, rows)

def at(self, x: float, y: float) -> float:
    i, fx = _bracket(self.x, x)
    j, fy = _bracket(self.y, y)
    a = self.rows[i][j] * (1 - fy) + self.rows[i][j + 1] * fy
    b = self.rows[i + 1][j] * (1 - fy) + self.rows[i + 1][j + 1] * fy
    return a * (1 - fx) + b * fx
\end{lstlisting}

\subsection{Aerodynamic drag and fuel flow}\label{subsec:route-flow}
\subsubsection*{Values}
Angle of attack, aerodynamic drag and required thrust at the current
weight and planned flight state. The model returns whole-aircraft fuel
flow and horizontal airspeed for the MFOB integration in
section~\ref{sec:mfob}.

\subsubsection*{Formulas}
Lift and drag use the coefficient definitions and steady straight-flight
force balance in \cite[pp.~1--1, 1--4 and 1--5]{how2004}.
Thrust is treated as parallel to the flight path. The engine-table
scaling and fuel-flow interpolation are app approximations.
\begin{align*}
 \alpha&=\alpha_i+
 \frac{C_L^*-C_{L,i}}{C_{L,i+1}-C_{L,i}}(\alpha_{i+1}-\alpha_i),\\
 D&=qS\,[C_D(M,\alpha)P_D(0,\alpha)+\Delta C_D],&
 T_{req}&=D+W\sin\gamma,\\
 Q&=Q_a+\max\!\left(0,\frac{T_{req}-T_a}{T_b-T_a}\right)(Q_b-Q_a),&
 V_h&=V_{TAS}\cos\gamma.
\end{align*}

\subsubsection*{Variables}
\begin{itemize}
\setlength{\itemsep}{0pt}
\setlength{\parsep}{0pt}
\interlinepenalty=10000
\item $h,M,q,V_{TAS}$: MSL altitude in feet, Mach number, dynamic
pressure in lbf/ft$^2$ and TAS in knots from
section~\ref{subsec:atmosphere}. Engine tables use $h$, not its metre
equivalent $z$.
\item $W,S$: current gross weight treated as weight in lbf, and the
aircraft's wing area in ft$^2$.
\item $s=\Delta h/\Delta d$, $\gamma=\arctan s$: leg gradient and
flight-path angle in radians; $\Delta h,\Delta d$ are the altitude
change and horizontal distance, both in feet.
\item $C_L^*=W\cos\gamma/(qS)$: lift coefficient required for force
balance normal to the flight path.
\item $C_L,C_D,P_L,P_D$: dimensionless base lift/drag coefficients and
their sideslip correction multipliers. Corrections use zero sideslip
and default to 1 when their table is absent.
\item $\alpha,\alpha_i$: computed angle of attack and sampled angles,
in degrees. Samples include the base lift-table breakpoints and
endpoints of its common range with the drag table, restricted to
$[-5^\circ,15^\circ]$.
\item $C_{L,i}=C_L(M,\alpha_i)P_L(0,\alpha_i)$: corrected sampled lift.
The samples must increase strictly; $i$ selects the interval enclosing
$C_L^*$, with $C_{L,i}\leq C_L^*\leq C_{L,i+1}$.
\item $DI,\Delta C_D$: takeoff displayed drag factor from
section~\ref{subsec:drag} and the added stores drag coefficient,
$\Delta C_D=(DI-1)k(M)$. The subtraction removes the clean-display baseline.
\item $k(M)$: coefficient increment per drag-index unit,
\[
 k(M)=k_0+\clip\!\left(\frac{M-M_0}{M_1-M_0},0,1\right)(k_1-k_0).
\]
$M_0,M_1,k_0,k_1$ are the Mach endpoints and their coefficients
(taken from BMS data);
$\clip$ clamps the fraction to $[0,1]$. A constant limiter has $k_0=k_1$.
\item $N_e,f_T,f_Q$: engine count and dimensionless thrust
and fuel-flow scale factors, fixed for the aircraft.
\item $A_T(h,\alpha)$: advanced thrust correction multiplier,
used when enabled; otherwise 1.
\item $r\in\{idle,MIL,AB\}$: engine regime.
$T_{table,r}(h,M),Q_{table,r}(h,M)$ are per-engine table values in lbf
and lb/h. Whole-aircraft regime values are
\[
 T_r=N_e f_T A_T(h,\alpha)T_{table,r}(h,M),\qquad
 Q_r=N_e f_Q Q_{table,r}(h,M).
\]
\item $a,b$: lower and upper interpolation regimes: idle/MIL when
$T_{req}\leq T_{MIL}$, otherwise MIL/AB. Below idle the model uses idle
flow; missing required AB data or thrust above AB makes it unavailable.
\item $D,T_{req}$: computed drag and required combined thrust, in lbf.
$Q,V_h$ are the returned combined fuel flow in lb/h and horizontal
airspeed in knots.
\end{itemize}

\subsubsection*{Data}
\begin{itemize}
\setlength{\itemsep}{0pt}
\setlength{\parsep}{0pt}
\item $h,M,q,V_{TAS}$ come from section~\ref{subsec:atmosphere}.
$s$ is computed from planned altitude changes and horizontal distances
in \path{lib/route_fuel.py:remaining_route_fuel}; the distances use the
first two coordinates of DTC \texttt{[STPT]} \texttt{target\_k} entries
in \path{<callsign/save>.ini}.
\item $W$ is computed from the selected CAM member's takeoff gross
weight and fuel, as described in section~\ref{subsec:gross-weight},
then adjusted for the current required fuel during the backward
integration in section~\ref{sec:mfob}. $DI$ stays at the saved takeoff
value from section~\ref{subsec:drag}.
\item $S$ is read from \path{airframe.area_sqft} in the aircraft's
\path{Sim/Acdata/*.txtpb}. $C_L,C_D$ are interpolated from
\path{aero_new2.clift_table} and \path{aero_new2.cdrag_table}, using
\texttt{mach\_bp}/\texttt{alpha\_bp} as fallback axes.
$P_L,P_D$ use \path{aero_new2.perc_lift_table} and
\path{aero_new2.perc_drag_table}, with
\texttt{beta\_perc\_bp}/\texttt{alpha\_bp} fallback axes.
All tables follow section~\ref{subsec:table-interpolation}.
\item $M_0,k_0,M_1,k_1$ are taken, in that order, from
\path{limits.limit.data.values} where \texttt{key=LIM\_STORES\_DRAG}
and \texttt{type=LIMITER\_LINE} in the same aircraft \texttt{.txtpb}.
For \texttt{LIMITER\_VALUE}, its single coefficient is used at both
endpoints. \path{lib/route_fuel.py:RoutePerformance.stores_cd} computes $\Delta C_D$.
\item $N_e,f_T,f_Q$ are read from \path{auxaero.engine.n_engines},
\path{engine.thrust_factor} and \path{engine.fuel_flow_factor}
in that file. $T_{table,r},Q_{table,r}$ use its
\path{engine.alt_mach_thrust_*_table} and
\path{engine.alt_mach_fuel_flow_*_table}, with \texttt{*} replaced by
\texttt{idle}, \texttt{mil} or \texttt{ab}. Missing axes inherit the
altitude/Mach axes from \path{engine.alt_mach_thrust_idle_table}.
\item $A_T$ is interpolated from
\path{engine.alt_alpha_thrust_factor_table} only when
\path{engine.has_advanced_thrust} is true in that \texttt{.txtpb};
its altitude axis falls back to the idle table's, and its angle axis
must be supplied by the table.
\item $\alpha,D,T_{req},Q,V_h$ are computed from these inputs in
\path{lib/route_fuel.py:RoutePerformance.flow}.
\texttt{from\_simulation} in that class requires
\texttt{fmtype=FM\_AFM} and \path{engine.has_fuel_flow} true in the
aircraft file. Out-of-range queries or invalid lift/engine ranges
make the route calculation unavailable.
\end{itemize}

\subsubsection*{Example}
Weasel7, flight 4/member 0 in \path{I/Campaign/weight_tests.cam}, has
takeoff weight 42689 lb and $DI=200$. Use that weight for this
illustrative level-flight sample ($s=\gamma=0$), with the STPT 3
atmosphere from section~\ref{subsec:atmosphere}.
\path{I/Sim/Acdata/F-16C-40 IAF.txtpb} gives $S=300$ ft$^2$ and
the limiter values $(0.9,0.0001,1.0,0.000283)$.
Since $M=0.527587<0.9$, $k(M)=0.0001$ and
$\Delta C_D=(200-1)(0.0001)=0.0199$.
The required lift and enclosing corrected table samples give
\[
 C_L^*=\frac{42689}{331.68072(300)}\simeq0.429017,\qquad
 \alpha=5+\frac{0.429017-0.3681381}{0.4443105-0.3681381}(6-5)
       \simeq5.7992^\circ.
\]
Here $P_L=P_D=A_T=1$; the drag table gives $C_D=0.04108642$, hence
\[
 D=T_{req}=331.68072(300)(0.04108642+0.0199)
          \simeq6068.406\ \text{lbf}.
\]
The aircraft has $N_e=f_T=f_Q=1$. Its scaled idle/MIL table values are
$(T_{idle},Q_{idle})=(-443.464,713.09)$ and
$(T_{MIL},Q_{MIL})=(13649.320,11741.231)$:
\[
 Q=713.09+\frac{6068.406+443.464}{13649.320+443.464}
       (11741.231-713.09)\simeq5808.88\ \text{lb/h},\qquad
 V_h=341.836\ \text{kt}.
\]
The actual MFOB integration reevaluates this model as required fuel
and therefore $W$ change.

\subsubsection*{Code}
In \texttt{from\_simulation}, \texttt{sim} is the parsed aircraft data
and \texttt{engines} is $N_e$. The returned object stores the wing
area, $N_ef_T$, $N_ef_Q$, interpolation tables and stores limiter.

\path{lib/route_fuel.py}, \texttt{RoutePerformance.from\_simulation}:
\begin{lstlisting}[numbers=left,firstnumber=80,numberstyle=\tiny,xleftmargin=2em]
return cls(
    _number(sim.airframe.wing_area_sqft, "wing area", positive=True),
    engines * _number(sim.engine.parameters.get("thrust_factor"), "thrust factor", positive=True),
    engines * _number(sim.engine.parameters.get("fuel_flow_factor"), "fuel-flow factor", positive=True),
    tables, values,
)
\end{lstlisting}

In \texttt{stores\_cd}, \texttt{index} is $DI-1$ and the limiter tuple
is $(M_0,k_0,M_1,k_1)$.
In \texttt{flow}, the inputs are altitude in feet, CAS in knots,
current weight, takeoff DI and leg gradient. \texttt{q\_area} is $qS$;
\texttt{alphas,lifts} hold $\alpha_i,C_{L,i}$.
The \texttt{\_bracket} call in \texttt{flow} locates the required lift
between samples. The local \texttt{engine} helper evaluates and scales a
regime table; \texttt{t0,t1,f0,f1} are $T_a,T_b,Q_a,Q_b$.
The returned pair is $(Q,V_h)$.

\path{lib/route_fuel.py}, \texttt{RoutePerformance.stores\_cd}:
\begin{lstlisting}[numbers=left,firstnumber=106,numberstyle=\tiny,xleftmargin=2em]
def stores_cd(self, mach: float, drag_index: float) -> float:
    index = _number(drag_index, "T/O drag factor") - 1
    if index < 0:
        raise FuelModelError("T/O drag factor")
    x0, y0, x1, y1 = self.stores_drag
    fraction = max(0., min(1., (mach - x0) / (x1 - x0)))
    return index * (y0 + fraction * (y1 - y0))
\end{lstlisting}

\path{lib/route_fuel.py}, \texttt{RoutePerformance.flow}:
\begin{lstlisting}[numbers=left,firstnumber=121,numberstyle=\tiny,xleftmargin=2em]
mach, tas, q = air_state(altitude_ft, cas_kt)
weight = _number(weight_lb, "gross weight", positive=True)
gamma = math.atan(gradient)
q_area = q * self.wing_area
cl, cd = self.tables["cl"], self.tables["cd"]
low, high = max(-5., cl.y[0], cd.y[0]), min(15., cl.y[-1], cd.y[-1])
alphas = sorted({low, high, *(a for a in cl.y if low < a < high)})

def correction(name, alpha):
    table = self.tables.get(name)
    return table.at(0, alpha) if table else 1.

lifts = tuple(cl.at(mach, a) * correction("perc_lift_table", a) for a in alphas)
if any(a >= b for a, b in zip(lifts, lifts[1:])):
    raise FuelModelError("Cruise lift range")
i, f = _bracket(lifts, weight * math.cos(gamma) / q_area)
alpha = alphas[i] + f * (alphas[i + 1] - alphas[i])
drag = (cd.at(mach, alpha) * correction("perc_drag_table", alpha)
        + self.stores_cd(mach, drag_index)) * q_area
required_thrust = drag + weight * math.sin(gamma)
factor = self.tables.get("alpha_thrust")
thrust_scale = self.thrust_scale * (factor.at(altitude_ft, alpha) if factor else 1.)

def engine(quantity, regime, scale):
    return self.tables[f"{quantity}_{regime}"].at(altitude_ft, mach) * scale

t0, t1 = (engine("thrust", r, thrust_scale) for r in ("idle", "mil"))
f0, f1 = (engine("fuel_flow", r, self.flow_scale) for r in ("idle", "mil"))
if t1 <= t0 or f0 <= 0 or f1 < f0:
    raise FuelModelError("Idle/MIL performance")
if required_thrust > t1:
    if not {"thrust_ab", "fuel_flow_ab"} <= self.tables.keys():
        raise FuelModelError("Afterburner performance")
    t0, f0 = t1, f1
    t1, f1 = engine("thrust", "ab", thrust_scale), engine("fuel_flow", "ab", self.flow_scale)
    if t1 <= t0 or f1 < f0 or required_thrust > t1:
        raise FuelModelError("Required thrust exceeds full AB")
fraction = max(0., (required_thrust - t0) / (t1 - t0))
return f0 + fraction * (f1 - f0), tas * math.cos(gamma)
\end{lstlisting}


\section{Liftoff speed (experimental)}\label{sec:takeoff}
\subsubsection*{Value}
The selected pilot's estimated liftoff speed, displayed as \texttt{L/O SPD}
in Admin, rounded to a whole knot. The tooltip specifies KCAS.
Unavailable values show a dash.

\subsubsection*{Formula}
The model assumes sea-level ISA, full afterburner, takeoff flaps and
horizontal airflow. The aircraft's takeoff pitch is treated as
angle of attack, with thrust along the body axis. Liftoff occurs when
lift plus vertical thrust supports takeoff gross weight:
\[
 C_{L,TO}(M,\alpha)=C_L(M,\alpha)P_L(0,\alpha)
   +C_{L,TEF}(M,\alpha)\frac{\delta_{TO}}{\delta_{max}},\]
\[ q(M_{LO})S C_{L,TO}(M_{LO},\alpha)
   +T(M_{LO})\sin\alpha=W_{TO},\]
\[V_{LO}=\frac{V_{ref}}{M_{ref}}M_{LO}.\]
The app finds the root by 60 bisection iterations over the Mach range
corresponding to 70--330 KCAS, after checking that support increases
strictly throughout that range.

\subsubsection*{Variables}
\begin{itemize}
\setlength{\itemsep}{0pt}
\setlength{\parsep}{0pt}
\interlinepenalty=10000
\item $V_{LO},M_{LO}$: computed liftoff speed in KCAS and its Mach number;
$M$ is a candidate Mach number.
\item $W_{TO},S$: takeoff gross weight from section~\ref{subsec:gross-weight}, treated as
weight in lbf, and the aircraft's wing area in ft$^2$.
\item $\alpha$: takeoff pitch in degrees (taken from BMS data), used as a
fixed angle of attack; the app requires $0<\alpha<30^\circ$.
\item $C_{L,TO},C_L,C_{L,TEF},P_L$: computed takeoff lift coefficient,
base lift coefficient, full trailing-edge-flap lift increment and
sideslip lift multiplier, respectively. All are dimensionless;
$P_L(0,\alpha)$ uses zero sideslip.
\item $\delta_{TO},\delta_{max}$: takeoff and maximum
trailing-edge-flap deflections in degrees. Their ratio is the applied
flap fraction, required to lie in $[0,1]$, with $\delta_{max}>0$.
\item $T(M)=N_e f_T T_{AB}(0,M)$: combined engine thrust in lbf.
$N_e,f_T$ are the engine count and thrust scale factor;
$T_{AB}$ is the per-engine full-afterburner thrust table, evaluated at
zero altitude in feet.
\item $V_{ref}=100$ KCAS: fixed reference speed chosen in the code;
$M_{ref},q_{ref}$ are its sea-level ISA Mach number and dynamic pressure,
computed using the atmosphere formulas in section~\ref{sec:flow}:
$M_{ref}\simeq0.15117647$, $q_{ref}\simeq33.85529$ lbf/ft$^2$.
\item $q(M)=q_{ref}(M/M_{ref})^2$: dynamic pressure in lbf/ft$^2$.
The reference scaling applies at the fixed sea-level ISA conditions.
\end{itemize}

\subsubsection*{Data}
\begin{itemize}
\setlength{\itemsep}{0pt}
\setlength{\parsep}{0pt}
\item $W_{TO}$ is computed from the saved member's fuel and loadout in
\path{<save>.cam} and the aircraft/store weights, using the files and
fields listed in section~\ref{subsec:gross-weight}. \path{lib/cam/loadout_data.py} passes
\texttt{gross\_weight\_lb} to the liftoff model.
\item $S,\alpha$ are read from \path{airframe.area_sqft} and
\path{auxaero.ai_behavior.opt_takeoff_pitch}, respectively, in the
aircraft's \path{Sim/Acdata/*.txtpb}. The parser
\path{lib/cam/opencam/aircraft_data.py} exposes them as
\texttt{wing\_area\_sqft} and \texttt{takeoff\_pitch\_deg}.
\item $\delta_{TO},\delta_{max}$ are read from
\path{auxaero.tef.tef_take_off} and \path{auxaero.tef.tef_max_angle},
respectively, in the same aircraft \texttt{.txtpb}.
\item $C_L,C_{L,TEF},P_L$ are interpolated from
\path{aero_new2.clift_table}, \path{aero_new2.c_lift_tef_table} and
\path{aero_new2.perc_lift_table}, respectively, in that \texttt{.txtpb}.
Their fallback axes within \texttt{aero\_new2} are
\texttt{mach\_bp}/\texttt{alpha\_bp},
\texttt{mach\_tef\_bp}/\texttt{alpha\_tef\_bp}, and
\texttt{beta\_perc\_bp}/\texttt{alpha\_bp}.
\path{lib/fuel.py:_Table} reads \path{data.values}, applies each table's
\texttt{multiplier} (1 if absent), and interpolates bilinearly; explicit
\path{x_axis.breakpoints}/\path{y_axis.breakpoints} take precedence.
\item $N_e,f_T$ are read from \path{auxaero.engine.n_engines} and
\path{engine.thrust_factor} in the aircraft \texttt{.txtpb}.
$T_{AB}$ is interpolated from \path{engine.alt_mach_thrust_ab_table},
using the same table rules. Missing altitude/Mach axes are taken from
\path{engine.alt_mach_thrust_idle_table} in that file.
\item $M_{ref},q_{ref}$ are computed by
\path{lib/route_fuel.py:air_state} at 0 ft and 100 KCAS, using fixed ISA
constants rather than saved weather or airfield elevation.
\path{lib/liftoff.py:LiftoffPerformance} computes $C_{L,TO}$, $T$ and
$V_{LO}$ from the preceding data. It requires \texttt{fmtype=FM\_AFM}
in the aircraft file, valid table coverage, increasing support and a
root within 70--330 KCAS. \path{lib/admin_brief.py} rounds the result.
\end{itemize}
Stores and CFTs enter through gross weight; DRAG is not an input.
The estimate omits ground effect, trim/CG, rotation dynamics, separate
store/CFT lift increments and advanced thrust corrections.

\subsubsection*{Example}
Falcon2-1 in \path{I/Campaign/weight_tests.cam} has
$W_{TO}=42256$ lb (section~\ref{subsec:gross-weight}).
Its \path{I/Sim/Acdata/F-16I-52+ IAF.txtpb} gives $S=300$ ft$^2$,
$\alpha=12^\circ$, $\delta_{TO}=\delta_{max}=21.5^\circ$, $N_e=1$ and
$f_T=1$. At the computed root $M_{LO}\simeq0.28798432$, the tables give
\[
 C_{L,TO}=0.8730039(1)+0.1239225(1)\simeq0.9969264.
\]
With $q\simeq122.85585$ lbf/ft$^2$ and $T\simeq26513.7865$ lbf,
\[
 122.85585(300)(0.9969264)+26513.7865\sin12^\circ
 \simeq42256\ \text{lbf}.
\]
Thus $V_{LO}=190.49546$ KCAS; Admin displays \texttt{190}.

For comparison, \path{B/Data/Sim/Acdata/F-15C.txtpb} gives $S=608$ ft$^2$,
$\alpha=10^\circ$, $\delta_{TO}=\delta_{max}=45^\circ$, $N_e=2$ and
$f_T=1$. With an illustrative $W_{TO}=50000$ lb, the model finds
$C_{L,TO}=0.7757$, $T=40982$ lbf and $V_{LO}=163.88281$ KCAS,
displayed as \texttt{164}. These are model estimates, not measured
liftoff speeds in BMS.

\subsubsection*{Code}
In \texttt{from\_simulation}, \texttt{sim} holds the parsed aircraft data;
\texttt{cl}, \texttt{flap} and \texttt{correction} are the three lift
tables, and \texttt{thrust} is the AB table.
\texttt{\_Table} implements the interpolation described above;
\texttt{\_number} validates numeric inputs.
The \texttt{mach} nodes combine the lift, flap and thrust Mach breakpoints
with the search limits. At each node, \texttt{lift} stores $C_{L,TO}$ and
\texttt{vertical} stores $T\sin\alpha$.
The resulting object stores $q_{ref}S/M_{ref}^2$ as
\texttt{q\_area\_per\_mach2}.

\path{lib/liftoff.py}, reference atmosphere:
\begin{lstlisting}[numbers=left,firstnumber=16,numberstyle=\tiny,xleftmargin=2em]
_REFERENCE_MACH, _, _REFERENCE_Q = air_state(0, 100)
_KCAS_PER_MACH = 100 / _REFERENCE_MACH
\end{lstlisting}

\path{lib/liftoff.py}, \texttt{LiftoffPerformance.from\_simulation}:
\begin{lstlisting}[numbers=left,firstnumber=31,numberstyle=\tiny,xleftmargin=2em]
area = _number(sim.airframe.wing_area_sqft, "wing area", positive=True)
alpha = _number(sim.takeoff_pitch_deg, "takeoff pitch", positive=True)
if alpha >= 30:
    raise ValueError("Takeoff pitch range")
tef = sim.flaps.get("tef", {})
flap_fraction = (_number(tef.get("tef_take_off"), "takeoff TEF")
                 / _number(tef.get("tef_max_angle"), "maximum TEF", positive=True))
if not 0 <= flap_fraction <= 1:
    raise ValueError("Takeoff TEF range")
aero = sim.aerodynamics.get("aero_new2")
if aero is None:
    raise ValueError("Liftoff lift tables required")
cl = _Table.read(aero.tables.get("clift_table"),
                 aero.breakpoints.get("mach_bp"), aero.breakpoints.get("alpha_bp"))
flap = _Table.read(aero.tables.get("c_lift_tef_table"),
                   aero.breakpoints.get("mach_tef_bp"), aero.breakpoints.get("alpha_tef_bp"))
correction = _Table.read(aero.tables.get("perc_lift_table"),
                         aero.breakpoints.get("beta_perc_bp"), aero.breakpoints.get("alpha_bp"))
lift_factor = correction.at(0, alpha)
idle = _Table.read(sim.engine.tables.get("alt_mach_thrust_idle_table"))
thrust = _Table.read(sim.engine.tables.get("alt_mach_thrust_ab_table"), idle.x, idle.y)
scale = (_number(sim.engine.number_of_engines, "engine count", positive=True)
         * _number(sim.engine.parameters.get("thrust_factor"), "thrust factor", positive=True)
         * math.sin(math.radians(alpha)))
low, high = (v / _KCAS_PER_MACH for v in (70., 330.))
mach = tuple(sorted({low, high, *(m for m in (*cl.x, *flap.x, *thrust.y) if low < m < high)}))
lift = tuple(_number(cl.at(m, alpha) * lift_factor + flap.at(m, alpha) * flap_fraction,
                     "takeoff lift coefficient", positive=True) for m in mach)
vertical = tuple(_number(thrust.at(0, m) * scale, "vertical AB thrust", positive=True) for m in mach)
result = cls(_REFERENCE_Q / _REFERENCE_MACH ** 2 * area, mach, lift, vertical)
\end{lstlisting}

The omitted validation checks that the derivative of support is positive
on every Mach interval. In \texttt{support}, \texttt{\_bracket} returns
the interval index and interpolation fraction; \texttt{thrust} there is
already the vertical component. The returned force is compared with
$W_{TO}$ in \texttt{speed\_kcas}, which converts the final Mach to KCAS.

\path{lib/liftoff.py}, \texttt{LiftoffPerformance.support/speed\_kcas}:
\begin{lstlisting}[numbers=left,firstnumber=78,numberstyle=\tiny,xleftmargin=2em]
def support(self, mach: float) -> float:
    i, fraction = _bracket(self.mach, mach)
    cl = self.lift[i] + fraction * (self.lift[i + 1] - self.lift[i])
    thrust = self.vertical_thrust[i] + fraction * (self.vertical_thrust[i + 1] - self.vertical_thrust[i])
    return self.q_area_per_mach2 * mach ** 2 * cl + thrust

def speed_kcas(self, gross_weight_lb: float) -> float:
    weight = _number(gross_weight_lb, "takeoff gross weight", positive=True)
    low, high = self.mach[0], self.mach[-1]
    if not self.support(low) <= weight <= self.support(high):
        raise ValueError("Liftoff speed outside 70-330 KCAS")
    for _ in range(60):
        middle = (low + high) / 2
        if self.support(middle) < weight:
            low = middle
        else:
            high = middle
    return (low + high) / 2 * _KCAS_PER_MACH
\end{lstlisting}


\section{Minimal fuel on board }\label{sec:mfob}
Minimal fuel on board (MFOB) is displayed in the flightplan. It is the minimal amount of fuel required to fly the rest of the flight plan from the current steerpoint + 500 lb. 
\subsubsection*{Value}
The selected pilot's required fuel on arrival at each steerpoint,
displayed as \texttt{MFOB} in the flightplan, in lb rounded up to the
next 100 lb. It includes any explicitly timed hold at that point.
Unavailable values and rows beyond the first planned landing show a dash.

\subsubsection*{Formulas}
Start with $F=F_{land}$ at the first planned landing and work backward.
At each steerpoint, add its hold requirement before recording $F_i$,
then integrate the preceding leg. Each step adds fuel using a midpoint
estimate of fuel weight:
\begin{align*}
 Q_1&=Q(h,V_{CAS},W_{dry}+F,DI,s),\\
 Q_2&=Q(h,V_{CAS},W_{dry}+F+Q_1\Delta t/2,DI,s),\\
 F_{new}&=F+Q_2\Delta t,\\
 \mathrm{MFOB}_i&=100\left\lceil\rnd(F_i/100,9)\right\rceil.
\end{align*}
Continue with the unrounded $F_{new}$. The landing reserve is included
once; the upward rounding is applied only to displayed cells.

\subsubsection*{Variables}
\begin{itemize}
\setlength{\itemsep}{0pt}
\setlength{\parsep}{0pt}
\interlinepenalty=10000
\item $i,L$: route-point index and first landing index, with
$i=1,\ldots,L$. For $i\geq2$, leg $i$ is flown from $i-1$ to $i$.
\item $F_i,\mathrm{MFOB}_i$: unrounded and displayed required fuel at
arrival at point $i$, in lb, including its hold and the landing reserve.
\item $F,F_{new}$: fuel requirement before and after one backward
integration step, in lb. Adding a step moves the estimate earlier
along the flight.
\item $F_{land}=500$ lb: fixed landing reserve.
The 100 lb display increment and nine decimal places in $\rnd$ are
fixed app settings. $\rnd(x,9)$ rounds $x$ to nine decimal places
before the ceiling $\lceil\cdot\rceil$, suppressing floating-point noise.
\item $W_{TO},F_{TO}$: the saved member's takeoff gross weight and fuel
from sections~\ref{subsec:gross-weight} and~\ref{subsec:takeoff-fuel}.
$W_{dry}=W_{TO}-F_{TO}$ is its fixed non-fuel weight, including retained
stores, tanks and CFT dry mass. All weights are in lb.
\item $DI$: fixed takeoff drag factor from section~\ref{subsec:drag}.
\item $N_i,E_i,Z_i$: DTC coordinates in feet; $N_i,E_i$ are north/east,
and $-Z_i$ supplies airport elevation. $h_i,V_{CAS,i}$ are the planned
MSL altitude in feet and CAS in knots assigned to point $i$.
\item $d_i$: horizontal leg length in NM, with
$c_{ft}=1852/0.3048$ feet per NM:
\[
 d_i=\frac{\sqrt{(N_i-N_{i-1})^2+(E_i-E_{i-1})^2}}{c_{ft}}.
\]
\item $n_i=\max(1,\lceil d_i/2\rceil)$: leg step count, giving
$\Delta d=d_i/n_i\leq2$ NM. $j=0,\ldots,n_i-1$ counts steps backward
from destination $i$.
\item $h,V_{CAS},s$: altitude in feet, CAS in knots and forward
flight-path gradient used for a leg step:
\[
 h=h_{i-1}+\left(1-\frac{j+1/2}{n_i}\right)(h_i-h_{i-1}),\qquad
 V_{CAS}=V_{CAS,i},\qquad
 s=\frac{h_i-h_{i-1}}{c_{ft}d_i}.
\]
Thus both flow evaluations use the same spatial midpoint. The gradient
retains the forward climb/descent sign during backward integration.
\item $Q(\cdot),Q_1,Q_2$: whole-aircraft fuel-flow model from
section~\ref{subsec:route-flow}, and its first and midpoint-weight
evaluations, in lb/h. $V_h$ is the horizontal airspeed in knots
returned with $Q_1$; a leg step uses $\Delta t=\Delta d/V_h$ hours.
\item $t_{arrival,i},t_{departure,i}$: arrival/departure clock times
in seconds after midnight. The hold duration in hours is
\[
 H_i=\frac{(t_{departure,i}-t_{arrival,i})\bmod86400}{3600}.
\]
The modulo allows a hold to cross midnight. Points without a timed
hold, including takeoff and landing, have $H_i=0$.
\item $n_{H,i}=\max(1,\lceil120H_i\rceil)$: step count for a nonzero
hold, giving $\Delta t=H_i/n_{H,i}\leq1/120$ h (30 s).
Hold steps use $h=h_i$, $V_{CAS}=V_{CAS,i}$ and $s=0$ in the same
fuel-weight update.
\end{itemize}

\subsubsection*{Data}
\begin{itemize}
\setlength{\itemsep}{0pt}
\setlength{\parsep}{0pt}
\item $W_{TO},F_{TO},DI$ come from the matched aircraft member in
\path{<save>.cam}, using the loadout data and files described in
section~\ref{sec:loadout}. \path{lib/fuel_brief.py:build_fuel_plan}
computes $W_{dry}$ from \texttt{gross\_weight\_lb} minus
\texttt{takeoff\_fuel\_lb}. Matching uses the briefing's package and
flight/callsign, then the selected pilot's roster position.
\item $N_i,E_i,Z_i,L$ come from the loaded DTC
\path{<callsign/save>.ini}, \texttt{[STPT]} entries
\texttt{target\_k=N,E,Z,action,...}, where $i=k+1$.
\path{lib/parsers/parse_callsign_ini.py} parses these entries;
\path{lib/fuel_brief.py:_route_points} stops at the lowest steerpoint
number with action 7 or 27. All points from 1 through $L$ are required.
\item $h_i,V_{CAS,i}$ come from the altitude and CAS columns of
\path{B/User/Briefings/briefing.txt}, parsed by
\path{lib/parsers/parse_briefing_txt.py}. An altitude suffix
\texttt{M} means thousands of feet; \texttt{A} (AGL) is unsupported.
Takeoff and landing use $h_i=-Z_i$ from the DTC instead.
A leg uses its destination row's CAS; takeoff borrows STPT 2's CAS,
and an absent landing row borrows the preceding row's.
\item $H_i$ is computed by \path{lib/fuel_brief.py:_hold_hours}
from the briefing row's arrival time and \texttt{Departure:} time in
a hold/departure comment. A patrol/attack label alone adds no time;
an explicit hold/departure comment with missing or invalid times
makes the calculation unavailable. Leg duration instead comes from
DTC distance and model airspeed.
\item $Q,V_h$ are computed from the flight state, current fuel weight
and takeoff $DI$ by \path{lib/route_fuel.py:RoutePerformance.flow},
using the aircraft's \path{Sim/Acdata/*.txtpb} fields described in
section~\ref{subsec:route-flow}.
\item $F_{land}$, the display increment and $c_{ft}$ are
\texttt{ALLOWANCE\_LB}, \texttt{ROUNDING\_LB} and \texttt{FEET\_PER\_NM}
in \path{lib/fuel.py}. The 2 NM/30 s step limits, nine-place rounding
and backward updates of $F$ are implemented in
\path{lib/route_fuel.py:remaining_route_fuel}.
A zero-length leg at unchanged altitude is skipped; a zero-length
altitude transition is rejected.
\item $\mathrm{MFOB}_i$ is attached to each briefing row by
\path{lib/fuel_brief.py:build_fuel_plan} and formatted with thousands
separators in \path{templates/flightplan.html}. Invalid route,
matching or performance data makes the planned values unavailable.
\end{itemize}
Store mass and $DI$ stay at takeoff values while fuel weight changes.
The estimate includes no wind, acceleration, bank penalty, refuelling
credit or store releases; altitude transitions are linear.

\subsubsection*{Example}
Use \path{I/Campaign/weight_tests.cam},
\path{I/Campaign/weight_tests.ini} and the 2026-10-12 11:57:32 export of
\path{B/User/Briefings/briefing.txt}: Weasel7, flight 4/member 0,
with \path{I/Sim/Acdata/F-16C-40 IAF.txtpb}.
The saved member gives
\[
 W_{dry}=42689-9202=33487\ \text{lb},\qquad DI=200.
\]
The first landing is STPT 10, where $F_{10}=500$ lb.
Working backward through the later legs reaches STPT 3's departure
with $F=6331.6719$ lb. Its arrival/departure times,
07:28:24--07:34:24, give
\[
 H_3=\frac{27264-26904}{3600}=0.1\ \text{h},\qquad
 n_{H,3}=\lceil120(0.1)\rceil=12,\qquad
 \Delta t=\frac{0.1}{12}=\frac1{120}\ \text{h}.
\]
At 5900 ft / 315 KCAS and $s=0$, the first backward hold step gives
$Q_1=5655.4490$ and $Q_2=5656.7086$ lb/h:
\[
 F_{new}=6331.6719+\frac{5656.7086}{120}
        \simeq6378.8112\ \text{lb}.
\]
After all 12 steps, the hold adds 567.0591 lb and
$F_3=6898.7310$ lb; the displayed MFOB is \texttt{6,900}.
Continuing backward gives:
\begin{center}\small
\begin{tabular}{@{}lrrrrrrrrrr@{}}\hline
STPT &1&2&3&4&5&6&7&8&9&10\\
MFOB (lb)&7,400&7,000&6,900&6,200&3,600&3,500&3,100&2,700&600&500\\\hline
\end{tabular}\end{center}
STPT 11 is the alternate landing and displays a dash.

\subsubsection*{Code}
In \texttt{build\_fuel\_plan}, \texttt{member} is the matched saved
aircraft, \texttt{model} is its \texttt{RoutePerformance} object, and
\texttt{route} contains the parsed \texttt{RoutePoint} records through
the first landing. The difference \texttt{gross - takeoff\_fuel} is
$W_{dry}$.

\path{lib/fuel_brief.py}, \texttt{build\_fuel\_plan}:
\begin{lstlisting}[numbers=left,firstnumber=189,numberstyle=\tiny,xleftmargin=2em]
gross = _planned_number(member.get("gross_weight_lb"), "T/O gross weight")
takeoff_fuel = _planned_number(member.get("takeoff_fuel_lb"), "T/O fuel")
if takeoff_fuel < 0:
    raise ValueError("T/O fuel")
route, landing_number = _route_points(rows, dtc)
values = remaining_route_fuel(model, route, gross - takeoff_fuel, member.get("drag_index"))
\end{lstlisting}

In \texttt{\_hold\_hours}, \texttt{row} is a parsed briefing steerpoint.
\texttt{arrival[1]} and \texttt{departure[1]} are the matched clock
strings; the local \texttt{seconds} helper converts them to seconds
after midnight.

\path{lib/fuel_brief.py}, \texttt{\_hold\_hours}:
\begin{lstlisting}[numbers=left,firstnumber=107,numberstyle=\tiny,xleftmargin=2em]
def _hold_hours(row: Any) -> float:
    comment = str(_attr(row, "comment", ""))
    # Timed patrol/attack repetition can also carry a Departure annotation.
    # Approximate that dwell at this point's planned speed and altitude.
    if not re.search(r"\b(?:hold|departure)\b", comment, re.IGNORECASE):
        return 0.
    departure = re.search(r"Departure:\s*(\d{2}:\d{2}:\d{2})z?", comment, re.IGNORECASE)
    arrival = re.fullmatch(r"(\d{2}:\d{2}:\d{2})z?", str(_attr(row, "time", "")).strip(), re.IGNORECASE)
    if not departure or not arrival:
        raise ValueError("Hold timing")

    def seconds(text):
        h, m, s = map(int, text.split(":"))
        if h > 23 or m > 59 or s > 59:
            raise ValueError("Hold timing")
        return h * 3600 + m * 60 + s

    return ((seconds(departure[1]) - seconds(arrival[1])) % 86400) / 3600
\end{lstlisting}

In \texttt{remaining\_route\_fuel}, \texttt{required} is $F$ and
\texttt{dry} is $W_{dry}$. \texttt{point} is destination $i$ and
\texttt{previous} is $i-1$; Python's \texttt{index} is zero-based.
\texttt{steps}, \texttt{step}, \texttt{ds} and \texttt{dt} correspond to
$n_i$, $j$, $\Delta d$ and $\Delta t$ for a leg; a hold uses
$n_{H,i}$ instead. \texttt{first,middle} are $Q_1,Q_2$.
\texttt{result[point.number]} records the rounded requirement after
adding the hold, while \texttt{required} retains full precision.

\path{lib/route_fuel.py}, \texttt{remaining\_route\_fuel}:
\begin{lstlisting}[numbers=left,firstnumber=180,numberstyle=\tiny,xleftmargin=2em]
dry = _number(dry_weight_lb, "dry weight", positive=True)
if not points:
    raise FuelModelError("Planned route")
required = float(ALLOWANCE_LB)
result = {}

def rounded():
    return ROUNDING_LB * math.ceil(round(required / ROUNDING_LB, 9))

for index in range(len(points) - 1, -1, -1):
    point = points[index]
    hold = _number(point.hold_hours, "hold duration")
    if hold < 0 or hold > 24:
        raise FuelModelError("Hold duration")
    if hold:
        steps = max(1, math.ceil(hold * 120))
        dt = hold / steps
        for _ in range(steps):
            first, _ = model.flow(point.altitude_ft, point.cas_kt, dry + required, drag_index)
            middle, _ = model.flow(point.altitude_ft, point.cas_kt,
                                   dry + required + first * dt / 2, drag_index)
            required += middle * dt
    result[point.number] = rounded()
    if index == 0:
        break
    previous = points[index - 1]
    distance = math.hypot(point.north_ft - previous.north_ft,
                          point.east_ft - previous.east_ft) / FEET_PER_NM
    if distance == 0:
        if point.altitude_ft != previous.altitude_ft:
            raise FuelModelError("Zero-length altitude transition")
        continue
    gradient = (point.altitude_ft - previous.altitude_ft) / (distance * FEET_PER_NM)
    steps = max(1, math.ceil(distance / 2))
    ds = distance / steps
    for step in range(steps):
        fraction = 1 - (step + .5) / steps
        altitude = previous.altitude_ft + fraction * (point.altitude_ft - previous.altitude_ft)
        first, speed = model.flow(altitude, point.cas_kt, dry + required, drag_index, gradient)
        dt = ds / speed
        middle, _ = model.flow(altitude, point.cas_kt, dry + required + first * dt / 2,
                               drag_index, gradient)
        required += middle * dt
return result
\end{lstlisting}

\section{Navigation, weather and terrain}\label{sec:nav}
\subsection{DTC offset geometry}\label{subsec:offset-geometry}
\subsubsection*{Values}
DTC offset-marker positions on the map, and the \texttt{BRG}/\texttt{RNG}
values updated when an offset is dragged: true bearing to $0.1^\circ$
and horizontal range to the nearest foot.

\subsubsection*{Formulas}
\begin{align*}
 N'&=N+r\cos\theta_T,& E'&=E+r\sin\theta_T,\\
 r&=\sqrt{\Delta N^2+\Delta E^2},&
 \theta_T&=\left(\frac{180^\circ}{\pi}\atan(\Delta E,\Delta N)\right)
             \bmod360^\circ.
\end{align*}
The first pair projects an offset; the second recovers bearing/range
from two positions. Geometry uses the theater's north/east plane.

\subsubsection*{Variables}
\begin{itemize}
\setlength{\itemsep}{0pt}
\setlength{\parsep}{0pt}
\interlinepenalty=10000
\item $N,E$ and $N',E'$: origin and offset north/east coordinates, in ft.
$\Delta N=N'-N$ and $\Delta E=E'-E$ are their differences.
\item $r$: horizontal range in ft; elevation is excluded.
\item $\theta_T$: clockwise bearing from true north, in degrees.
$\atan(\Delta E,\Delta N)$ returns radians; the explicit factor
$180^\circ/\pi$ converts to degrees. Projection converts $\theta_T$
back to radians before evaluating sine/cosine.
\end{itemize}

\subsubsection*{Data}
\begin{itemize}
\setlength{\itemsep}{0pt}
\setlength{\parsep}{0pt}
\item $N,E$ are read from the first two fields of
\texttt{[STPT].target\_k=N,E,Z,action,...} in the loaded
\path{<callsign/save>.ini}; the displayed steerpoint number is $k+1$.
The parser is \path{lib/parsers/parse_callsign_ini.py}.
\item $\theta_T,r$ come from \texttt{[NAV OFFSETS]} in the same file,
or the edited DTC fields. \texttt{VIP}, \texttt{VIPPUP}, \texttt{VRP}
and \texttt{VRPPUP} rows store \texttt{stpt,brg,rng,elev};
\texttt{OA1-stpt}/\texttt{OA2-stpt} store \texttt{brg,rng,elev}.
Here \texttt{brg} is true degrees and \texttt{rng} is feet.
\item $N',E'$ are computed by \texttt{projectOffsetFromBmsCoord}
or \texttt{projectBearingRange} in \path{templates/dtc.html}.
Dragging supplies the new point through \texttt{mapLatLngToBmsCoord};
\texttt{getDraggedOffsetMetrics} applies the inverse formulas.
Pop-up geometry in section~\ref{sec:popup} uses the same projection
and \texttt{bearingRangeBetween}.
\item The conversion between $N,E$ and web-map latitude/longitude
uses the theater's WGS84 Transverse Mercator projection. Its construction
follows oakdesign's \emph{Falcon BMS Tacview Converter}~\cite{oakdesign_converter},
as credited in \path{templates/map.html} (line 646). The DTC reverse
conversion is \texttt{mapLatLngToBmsCoord} in \path{templates/dtc.html}.
\end{itemize}

\subsubsection*{Example}
Using STPT 1 and 2 from \path{I/Campaign/weight_tests.ini} as the two
positions (these are route points, not saved offset rows):
\[
 (N,E)=(1923708.25,1733469.375),\qquad
 (N',E')=(1999147.75,1687549.75)\ {\rm ft}.
\]
Thus
\[
 r=\sqrt{75439.5^2+(-45919.625)^2}=88316.08076\ {\rm ft},\qquad
 \theta_T=328.67135^\circ.
\]
The offset formatter gives \texttt{BRG 328.7}, \texttt{RNG 88316}.
Projecting the unrounded $r,\theta_T$ from STPT 1 recovers STPT 2.

\subsubsection*{Code}
\path{templates/dtc.html}, \texttt{projectBearingRange} and
\texttt{bearingRangeBetween}; \texttt{normalizeBearing} wraps
degrees into $[0,360)$:
\begin{lstlisting}[numbers=left,firstnumber=742,numberstyle=\tiny,xleftmargin=2em]
function projectBearingRange(baseCoord, bearing, rangeFeet) {
  if (!baseCoord || !Number.isFinite(bearing) || !Number.isFinite(rangeFeet)) return null;
  const radians = normalizeBearing(bearing) * Math.PI / 180;
  return {
    northFeet: baseCoord.northFeet + Math.cos(radians) * rangeFeet,
    eastFeet: baseCoord.eastFeet + Math.sin(radians) * rangeFeet,
  };
}

function bearingRangeBetween(fromCoord, toCoord) {
  if (!fromCoord || !toCoord) return null;
  const deltaNorth = toCoord.northFeet - fromCoord.northFeet;
  const deltaEast = toCoord.eastFeet - fromCoord.eastFeet;
  const rangeFeet = Math.hypot(deltaNorth, deltaEast);
  let bearing = Math.atan2(deltaEast, deltaNorth) * 180 / Math.PI;
  if (bearing < 0) bearing += 360;
  return {
    brg: formatDraggedBearing(bearing),
    rng: formatDraggedRange(rangeFeet),
  };
}
\end{lstlisting}

The same file's formatters set the precision returned in
\texttt{brg}/\texttt{rng} and used for dragged offsets:
\begin{lstlisting}[numbers=left,firstnumber=923,numberstyle=\tiny,xleftmargin=2em]
function formatDraggedBearing(value) {
  if (!Number.isFinite(value)) return "";
  const rounded = Math.round(value * 10) / 10;
  return Number.isInteger(rounded) ? String(rounded) : rounded.toFixed(1);
}

function formatDraggedRange(value) {
  if (!Number.isFinite(value)) return "";
  return String(Math.round(value));
}
\end{lstlisting}

\subsection{Magnetic headings}\label{subsec:magnetic-headings}
\subsubsection*{Values}
Flight-plan \texttt{MHDG}, in whole degrees, and pop-up magnetic
headings, to $0.1^\circ$. When variation is unavailable, the app
identifies the headings as true (\texttt{THDG} or a $^\circ$T suffix).

\subsubsection*{Formulas}
\begin{align*}
 \theta_M&=(\theta_T-\nu)\bmod360^\circ,\\
 \theta_{FP}&=\lfloor\theta_M+0.5^\circ\rfloor\bmod360^\circ,\\
 \theta_{POP}&=\left(\frac{\lfloor10\theta_M+0.5^\circ\rfloor}{10}\right)
                  \bmod360^\circ.
\end{align*}
Rounding operates on the numerical degree values.
For a pilot-entered magnetic heading, the inverse conversion
$\theta_T=(\theta_M+\nu)\bmod360^\circ$ is applied before geometry.

\subsubsection*{Variables}
\begin{itemize}
\setlength{\itemsep}{0pt}
\setlength{\parsep}{0pt}
\interlinepenalty=10000
\item $\theta_T,\theta_M$: true and magnetic heading, in degrees.
$\theta_{FP},\theta_{POP}$ are the rounded flight-plan and pop-up
headings.
\item $\nu$: magnetic variation in degrees, positive east and negative
west. It is bilinearly interpolated:
\[
 \nu=(1-a)(1-b)\nu_{00}+a(1-b)\nu_{10}
       +(1-a)b\nu_{01}+ab\nu_{11}.
\]
\item $x=N/c_k,\ y=E/c_k$: the magnetic grid's coordinates in km,
where $N,E$ are DTC north/east feet and
$c_k=1000/0.3048=3280.839895013123$ ft/km.
Here the grid's $x$ is \emph{north}, and $y$ is \emph{east}.
\item $x_0,x_1$ and $y_0,y_1$: adjacent grid-axis entries bracketing
$x,y$. The fractional positions are
$a=(x-x_0)/(x_1-x_0)$ and $b=(y-y_0)/(y_1-y_0)$.
\item $\nu_{00},\nu_{10},\nu_{01},\nu_{11}$: grid variations at
$(x_0,y_0)$, $(x_1,y_0)$, $(x_0,y_1)$, $(x_1,y_1)$.
\end{itemize}

\subsubsection*{Data}
\begin{itemize}
\setlength{\itemsep}{0pt}
\setlength{\parsep}{0pt}
\item $\theta_T$ for a flight-plan row is the heading column in
\path{B/User/Briefings/briefing.txt}, parsed by
\path{lib/parsers/parse_briefing_txt.py}. It is not recomputed from
DTC positions. Pop-up headings come from the geometry described in
section~\ref{sec:popup}.
\item $N,E$ come from the matching numbered DTC steerpoint in
\path{<callsign/save>.ini}. \path{lib/theater_paths.py:magnetic_variation_for_steerpoints}
samples variation at each flight-plan row's point. The pop-up plan
uses the selected reference steerpoint's variation throughout.
\item $x_0,x_1,y_0,y_1$ and the four corner variations are read from
the active terrain's \path{Weather/MagVarMap_*.CSV} by
\path{openchart/vendor/opencam/theater.py:load_magnetic_variation_grid}.
The header after \texttt{y\textbackslash x} gives the $x$ axis;
the first entry of each subsequent row gives $y$; the remaining
entries are signed degrees. The leading \texttt{\#year} is retained
as the grid epoch; no date-dependent correction is applied.
\item $\nu$ is computed by \texttt{MagneticVariationGrid} in that
file. A missing grid or an out-of-grid point leaves variation
unavailable; the app then retains true headings.
\end{itemize}

\subsubsection*{Example}
At STPT 2 of \path{I/Campaign/weight_tests.ini},
$(x,y)=(609.3402342,514.3651638)$ km.
\path{I/Terrdata/Israel/Weather/MagVarMap_Israel.CSV} brackets these
by $(x_0,x_1)=(608,624)$ and $(y_0,y_1)=(512,528)$ km, with
\[
 (\nu_{00},\nu_{10},\nu_{01},\nu_{11})
 =(4.145995,4.144043,4.162748,4.160932)^\circ.
\]
Therefore $a=0.08376464$, $b=0.14782274$, and
$\nu=4.14830965^\circ$. The briefing's STPT 2 heading is $329^\circ$:
\[
 \theta_M=329-4.14830965=324.85169035^\circ
 \quad\Longrightarrow\quad \texttt{MHDG 325}.
\]

\subsubsection*{Code}
\path{openchart/vendor/opencam/theater.py},
\texttt{MagneticVariationGrid}: \texttt{x\_fraction}/\texttt{y\_fraction}
are $a,b$; \texttt{lower}/\texttt{upper} interpolate the two $y$ rows.
The campaign-coordinate wrapper shows the north/east axis order.
\begin{lstlisting}[numbers=left,firstnumber=154,numberstyle=\tiny,xleftmargin=2em]
def variation_degrees(self, x_km: float, y_km: float) -> float:
    """Bilinearly interpolate signed magnetic variation on the grid."""

    x = _bounded_axis_value(x_km, self.x_km, "x_km")
    y = _bounded_axis_value(y_km, self.y_km, "y_km")
    x_index, x_fraction = _interpolation_position(x, self.x_km)
    y_index, y_fraction = _interpolation_position(y, self.y_km)
    lower_left = self.values_degrees[y_index][x_index]
    lower_right = self.values_degrees[y_index][x_index + 1]
    upper_left = self.values_degrees[y_index + 1][x_index]
    upper_right = self.values_degrees[y_index + 1][x_index + 1]
    lower = lower_left + (lower_right - lower_left) * x_fraction
    upper = upper_left + (upper_right - upper_left) * x_fraction
    return lower + (upper - lower) * y_fraction

def variation_at_campaign_position(
    self,
    position_x_ft: float,
    position_y_ft: float,
) -> float:
    return self.variation_degrees(
        position_x_ft / FEET_PER_KILOMETER,
        position_y_ft / FEET_PER_KILOMETER,
    )
\end{lstlisting}

\path{lib/html_gen.py}, \texttt{\_flightplan\_headings}, the branch
where variation is available:
\begin{lstlisting}[numbers=left,firstnumber=37,numberstyle=\tiny,xleftmargin=2em]
    heading = str(math.floor((true_heading - variation) % 360 + 0.5) % 360)
headings[str(point.number)] = heading
\end{lstlisting}

\path{templates/dtc_popup.js}, \texttt{formatPopupHeading};
a missing variation contributes zero and is labelled true by
\texttt{popupHeadingReference}:
\begin{lstlisting}[numbers=left,firstnumber=12,numberstyle=\tiny,xleftmargin=2em]
function formatPopupHeading(trueHeading, plan) {
  const heading = Number.parseFloat(trueHeading);
  if (!Number.isFinite(heading)) return "---";
  const magnetic = normalizeBearing(heading - (plan.magneticVariation ?? 0));
  return formatPopupDegrees((Math.round(magnetic * 10) / 10) % 360);
}
\end{lstlisting}

The inverse conversion in the same file preserves true bearings
for geometry and stored DTC offsets:
\begin{lstlisting}[numbers=left,firstnumber=441,numberstyle=\tiny,xleftmargin=2em]
const magneticVariation = getPopupMagneticVariation();
// Convert the pilot's heading once; geometry and DTC offset rows stay true.
const trueHeading = normalizeBearing(heading + (magneticVariation ?? 0));
\end{lstlisting}

\subsection{Departure weather}\label{subsec:departure-weather}
\subsubsection*{Values}
The ATIS table displays \texttt{WIND} (whole degrees/knots),
\texttt{TEMP} (whole degrees Celsius), \texttt{QNH}
(two decimals inHg and whole hPa), \texttt{WX} (category),
\texttt{LOW CLD BASE} (whole ft MSL or \texttt{NONE}), and
\texttt{FOG BOUNDARY} (km to one decimal).

\subsubsection*{Formulas}
\[
 QNH_{\rm inHg}=\frac{p}{C_p},\qquad
 \theta_W=\lfloor(\theta\bmod360^\circ)+0.5^\circ\rfloor
               \bmod360^\circ.
\]
Pressure is converted directly from the selected cell; the app
does not apply an airport-elevation correction. Wind direction
is wrapped and rounded without magnetic-variation correction.
Other weather values are formatted from the stored cell.

\subsubsection*{Variables}
\begin{itemize}
\setlength{\itemsep}{0pt}
\setlength{\parsep}{0pt}
\interlinepenalty=10000
\item $p,QNH_{\rm inHg}$: cell pressure in hPa and its inHg equivalent.
$C_p=33.8638866667$ hPa/inHg is the application's fixed conversion
constant, \texttt{HPA\_PER\_INHG}.
\item $\theta,\theta_W,U$: stored first-layer wind direction,
rounded displayed direction (degrees), and first-layer speed (knots).
\item $T,\kappa,H_c,F$: cell temperature ($^\circ$C), weather-type
code, low-cloud base (ft MSL), and stored fog-end distance (km).
$F$ is shown as a fog boundary, not a computed visibility.
\item $N,E$: departure north/east coordinates in ft;
$N_k=0.0003048N$, $E_k=0.0003048E$ convert them to km.
\item $L$: theater side length in km; $w,h$: FMAP width and height
in cells.
\item $c,r$: zero-based column and row of the sampled cell:
\[
 c=\left\lfloor wE_k/L-0.5\right\rfloor\bmod w,\qquad
 r=\left\lceil h(1-N_k/L)-0.5\right\rceil\bmod h.
\]
Columns run eastward and rows southward. The fixed half-cell
offset and edge wrapping implement the stored-cell convention
used by \texttt{\_departure\_cell}; coordinates must be within
$0\leq N_k,E_k\leq L$.
\end{itemize}

\subsubsection*{Data}
\begin{itemize}
\setlength{\itemsep}{0pt}
\setlength{\parsep}{0pt}
\item $N,E$ are read from the unique steerpoint with action 1 in
\texttt{[STPT]} of the loaded DTC \path{<callsign/save>.ini}.
$L$ comes from the active theater's
\texttt{Theater size in KM} in \path{NewTerrain/Theater.txt},
through \path{lib/theater_paths.py:read_theater_map_info}.
\item $w,h,p,T,\theta,U,\kappa,H_c,F$ come from the loaded CAM's
same-named adjacent \path{<save>.fmap}, parsed by
\path{lib/cam/opencam/fmap_parser.py} and
\path{lib/cam/opencam/fmap_wrappers.py}.
The selected cell has array index $rw+c$. Its fields are
\texttt{pressure\_hpa}, \texttt{temperature\_c},
\texttt{wind\_directions\_deg[0]}, \texttt{wind\_speeds\_knots[0]},
\texttt{weather\_type}, \texttt{cumulus\_base\_ft} and
\texttt{fog\_end\_km}. No spatial interpolation or movement in time
is applied to the saved snapshot.
\item $\kappa$ is mapped in \path{lib/atis.py} as
1 sunny, 2 fair, 3 poor, 4 inclement. Sunny displays
\texttt{NONE} for low cloud; types 2--4 display a valid $H_c$.
Unrecognised codes display \texttt{UNKNOWN (code)}.
\item $QNH_{\rm inHg},\theta_W$ and the formatted cell values are
computed by \path{lib/atis.py:build_departure_atis} and displayed
by \path{templates/atis.html}. Missing or invalid fields remain
unavailable; a missing/mismatched FMAP produces a status message.
\end{itemize}

\subsubsection*{Example}
\path{B/Data/Campaign/714_html_test.fmap} has $w=h=59$.
For the departure in \path{test/User/Config/wsy.ini},
$(N,E)=(1700669.625,2468185)$ ft, and Korea's
\path{B/Data/TerrData/Korea/NewTerrain/Theater.txt} gives $L=1024$ km:
\begin{align*}
 (N_k,E_k)&=(518.3641017,752.302788)\ {\rm km},\\
 c&=\lfloor59(752.302788)/1024-0.5\rfloor\bmod59=42,\\
 r&=\lceil59(1-518.3641017/1024)-0.5\rceil\bmod59=29.
\end{align*}
Cell $29\cdot59+42=1753$ has
$p=1011.96533203$ hPa, $\theta=39.63571167^\circ$,
$U=9.66757011$ kt, $T=5.97265148^\circ$C, $\kappa=3$,
$H_c=2159.93506$ ft and $F=59.71747$ km. In particular,
\[
 QNH_{\rm inHg}=1011.96533203/33.8638866667
              =29.88332\longrightarrow29.88.
\]
The readout is \texttt{WIND: 040 DEG / 10 KT},
\texttt{TEMP: 6 DEG C}, \texttt{QNH: 29.88 INHG / 1012 HPA},
\texttt{WX: POOR}, \texttt{LOW CLD BASE: 2,160 FT MSL},
\texttt{FOG BOUNDARY: 59.7 KM}.

\subsubsection*{Code}
\path{lib/atis.py}, \texttt{\_departure\_cell}, coordinate conversion
and cell selection:
\begin{lstlisting}[numbers=left,firstnumber=34,numberstyle=\tiny,xleftmargin=2em]
north_km, east_km = north_ft * 0.0003048, east_ft * 0.0003048
if not (0 <= north_km <= theater_size_km and 0 <= east_km <= theater_size_km):
    raise ValueError("Departure is outside the active theater")
column = math.floor(east_km / theater_size_km * weather.width - 0.5) % weather.width
row = math.ceil((1 - north_km / theater_size_km) * weather.height - 0.5) % weather.height
return weather.cell(row, column)
\end{lstlisting}

The same file's \texttt{build\_departure\_atis} converts and formats
the fields from that cell; \texttt{\_finite} rejects non-finite values:
\begin{lstlisting}[numbers=left,firstnumber=84,numberstyle=\tiny,xleftmargin=2em]
result["weather_type"] = WEATHER_TYPES.get(cell.weather_type, f"Unknown ({cell.weather_type})")
# The per-cell cloud-base field controls the low cloud layer (ft MSL).
# Sunny disables that layer even when the FMAP retains a base value.
cloud_base = _finite(cell.cumulus_base_ft)
if cell.weather_type == 1:
    result["cloud_base"] = "None"
elif cell.weather_type in (2, 3, 4) and cloud_base is not None and cloud_base >= 0:
    result["cloud_base"] = f"{cloud_base:,.0f} ft MSL"
# Preserve the stored fog-end distance. Live ATIS visibility has not
# been established as a direct conversion of this input to BMS's model.
fog_boundary = _finite(cell.fog_end_km)
if fog_boundary is not None and fog_boundary >= 0:
    result["fog_boundary"] = f"{fog_boundary:.1f} km"

pressure, temperature = _finite(cell.pressure_hpa), _finite(cell.temperature_c)
speed, direction = _finite(cell.wind_speeds_knots[0]), _finite(cell.wind_directions_deg[0])
if speed is not None and speed >= 0 and direction is not None:
    result["wind"] = f"{math.floor(direction % 360 + 0.5) % 360:03d} deg / {speed:.0f} kt"
if temperature is not None:
    result["temperature"] = f"{temperature:.0f} deg C"
if pressure is not None and pressure > 0:
    result["qnh"] = f"{pressure / HPA_PER_INHG:.2f} inHg / {pressure:.0f} hPa"
\end{lstlisting}

\subsection{Terrain and AGL-to-MSL conversion}\label{subsec:terrain}
\subsubsection*{Values}
The pop-up calculator displays AGL and MSL altitudes in whole feet
for \texttt{PDWN}, \texttt{APEX}, \texttt{RALT} and \texttt{MRA};
map labels also use this conversion.

\subsubsection*{Formula}
\[
 H_{\rm MSL}=H_{\rm AGL}+H_{\rm terrain}.
\]

\subsubsection*{Variables}
\begin{itemize}
\setlength{\itemsep}{0pt}
\setlength{\parsep}{0pt}
\interlinepenalty=10000
\item $H_{\rm AGL},H_{\rm MSL}$: height above ground and altitude
above mean sea level, in ft.
\item $H_{\rm terrain}$: signed ground elevation in ft, read from
one heightmap sample at the requested position.
\item $N,E$: requested north/east coordinates in ft.
$N_{\max}=E_{\max}=1000L/0.3048$ are the theater extents in ft,
where $L$ is the theater side length in km.
\item $w,h$: heightmap width and height in pixels.
$x,y$: zero-based pixel column and row, computed as
\begin{align*}
 x&=\clip\!\left(\trunc\!\left[(w-1)E/E_{\max}\right],0,w-1\right),\\
 y&=\clip\!\left(\trunc\!\left[(h-1)(1-N/N_{\max})\right],0,h-1\right).
\end{align*}
$\trunc$ truncates to an integer; $\clip$ clamps to the pixel bounds.
Columns run eastward and rows southward. Out-of-theater positions
are rejected before this calculation.
\item $o=2(yw+x)$: byte offset of the sample in the RAW file.
The factor 2 is the byte size of one signed 16-bit elevation.
\end{itemize}

\subsubsection*{Data}
\begin{itemize}
\setlength{\itemsep}{0pt}
\setlength{\parsep}{0pt}
\item $H_{\rm AGL},N,E$ come from the pop-up plan in
\path{templates/dtc_popup.js}; its inputs and geometry are described
in section~\ref{sec:popup}. \texttt{PDWN}/\texttt{APEX} use terrain
at the pull-down point; \texttt{RALT}/\texttt{MRA} use terrain at
the release point (or target if that point is absent).
\item $L,w,h$ are read from \texttt{Theater size in KM} and
\texttt{Map size in pixels} in the active
\path{NewTerrain/Theater.txt}, resolved through the theater TDF's
\texttt{terraindir}. \path{lib/terrain_elevation.py:_resolve_heightmap_spec}
sets both extents from $L$ and both dimensions from the pixel size.
If the pixel size is absent, it infers a square side from half the
RAW byte count, falling back to 32768 if that is not a square.
\item $H_{\rm terrain}$ is read from the active
\path{NewTerrain/Heightmaps/Heightmap.RAW} by
\path{lib/terrain_elevation.py:HeightmapSampler.elevation_feet}.
The two bytes at $o$ are decoded as a little-endian signed 16-bit
integer in feet. This is a single-pixel lookup without interpolation.
The heightmap lookup approach is also credited to oakdesign's
converter~\cite{oakdesign_converter} in \path{README.md};
see its \texttt{CoordinateConverter.get\_elevation}.
\item $H_{\rm MSL}$ is computed by \texttt{popupAglMslLine} and
related formatters in \path{templates/dtc_popup.js}. Until a finite
terrain sample is available, the label reads \texttt{MSL pending}.
\end{itemize}

\subsubsection*{Example}
\path{I/Terrdata/Israel/NewTerrain/Theater.txt} gives $L=1024$ km
and $w=h=32768$, hence $N_{\max}=E_{\max}=3359580.05249$ ft.
At STPT 1 of \path{I/Campaign/weight_tests.ini},
$(N,E)=(1923708.25,1733469.375)$ ft:
\[
 x=16907,\qquad y=14004,\qquad
 o=2(14004\cdot32768+16907)=917799958.
\]
The sample in
\path{I/Terrdata/Israel/NewTerrain/Heightmaps/Heightmap.RAW}
is $H_{\rm terrain}=154$ ft. An illustrative $H_{\rm AGL}=2000$ ft
at that position would display $H_{\rm MSL}=2000+154=2154$ ft.
The airport elevation in the DTC is instead $-Z=120$ ft;
section~\ref{sec:mfob} uses that DTC value.

\subsubsection*{Code}
\path{lib/terrain_elevation.py}, \texttt{HeightmapSampler.elevation\_feet},
after the theater-bound check: map to a pixel, seek to its two bytes,
and decode the signed elevation.
\begin{lstlisting}[numbers=left,firstnumber=59,numberstyle=\tiny,xleftmargin=2em]
east_ratio = east_feet / self.spec.max_east_feet if self.spec.max_east_feet else 0
north_ratio = north_feet / self.spec.max_north_feet if self.spec.max_north_feet else 0
pixel_x = int(east_ratio * (self.spec.width - 1))
pixel_y = int((1.0 - north_ratio) * (self.spec.height - 1))
pixel_x = max(0, min(self.spec.width - 1, pixel_x))
pixel_y = max(0, min(self.spec.height - 1, pixel_y))

offset = (pixel_y * self.spec.width + pixel_x) * 2
with self._lock:
    if self._file is None:
        try:
            self._file = self.spec.path.open("rb")
        except OSError as exc:
            raise HeightmapReadError(f"Could not open heightmap {self.spec.path}: {exc}") from exc
    try:
        self._file.seek(offset)
        raw = self._file.read(2)
    except OSError as exc:
        raise HeightmapReadError(f"Could not read heightmap {self.spec.path}: {exc}") from exc

if len(raw) != 2:
    raise HeightmapReadError(f"Expected 2 bytes at offset {offset}, got {len(raw)}")
return int(struct.unpack("<h", raw)[0])
\end{lstlisting}

\path{templates/dtc_popup.js}, \texttt{popupAglMslLine};
\texttt{aglFeet} and \texttt{terrainFeet} are $H_{\rm AGL}$ and
$H_{\rm terrain}$:
\begin{lstlisting}[numbers=left,firstnumber=252,numberstyle=\tiny,xleftmargin=2em]
function popupAglMslLine(label, aglFeet, terrainFeet) {
  const aglText = formatPopupNumber(aglFeet, 0) + " ft";
  const mslText = Number.isFinite(terrainFeet)
    ? formatPopupNumber(aglFeet + terrainFeet, 0) + " ft"
    : "MSL pending";
  return popupLine(label, popupMetric("AGL", aglText), popupMetric("MSL", mslText));
}
\end{lstlisting}


\section{Pop-up attack calculator}\label{sec:popup}
The calculator uses pilot-entered attack parameters and fixed application
constants. BMS data supplies the reference positions, magnetic variation
and terrain elevations; the model does not read weapon or aircraft
aerodynamic tables.

\subsection{Bomb range}\label{subsec:popup-bomb}
\subsubsection*{Value}
\texttt{BOMB RANGE} in the DTC pop-up inputs and support values, in whole
feet. Low/high-drag modes compute it; manual mode uses the entered range.

\subsubsection*{Formulas}
For each time step before ground crossing:
\begin{align*}
 a_j&=\begin{cases}
 -kA_d|u_j|/\sqrt{u_j^2+\epsilon},&u_j>0,\\
 0,&u_j=0,
 \end{cases}\\
 \Delta x_j&=u_j\Delta t+\tfrac12a_j\Delta t^2,&
 \Delta h_j&=-w_j\Delta t-\tfrac12g\Delta t^2,\\
 x_{j+1}&=x_j+\Delta x_j,&h_{j+1}&=h_j+\Delta h_j,\\
 u_{j+1}&=\max(0,u_j+a_j\Delta t),&
 w_{j+1}&=w_j+g\Delta t.
\end{align*}
The final step is corrected by its overshoot fraction:
\[
 B=x_m-f\Delta x_{m-1},\qquad t_f=(m-f)\Delta t.
\]

\subsubsection*{Variables}
\begin{itemize}
\setlength{\itemsep}{0pt}
\setlength{\parsep}{0pt}
\interlinepenalty=10000
\item $H,V,\delta$: release height in ft AGL, release TAS in knots,
and positive dive angle in degrees.
\item $x_j,h_j,u_j,w_j$: accumulated horizontal range, remaining
height (ft), horizontal speed and downward speed (ft/s) at step $j$.
Initially $x_0=0$, $h_0=H$, and
\[
 u_0=\rnd_{0.01}(c_VV\cos\delta),\qquad
 w_0=\rnd_{0.01}(c_VV\sin\delta).
\]
$\rnd_{0.01}$ rounds to $0.01$ ft/s;
$c_V=1.68781$ ft/s per knot is \texttt{FEET\_PER\_KNOT\_SECOND}.
\item $\Delta t=0.1$ s: fixed time step. $j=0,\ldots,m-1$ is the
step index; $m$ is the first index for which $h_m\leq0$.
\item $a_j$: horizontal acceleration in ft/s$^2$;
$\Delta x_j,\Delta h_j$: range and signed height changes during step $j$.
\item $k,g$: fixed model parameters, $(0.165,32.177)$ for low drag
and $(1,20.91505)$ for high drag. $k$ is dimensionless and $g$ is
the model's downward acceleration in ft/s$^2$; the high-drag value
is an effective acceleration.
$A_d=140$ ft/s$^2$ and $\epsilon=0.1$ (ft/s)$^2$ are fixed constants
in the horizontal-drag expression.
\item $f=h_m/\Delta h_{m-1}$: fraction of the final step below
ground. Both numerator and denominator are nonpositive.
$B$ is the corrected range in ft and $t_f$ the corrected fall time
in seconds. The latter feeds \texttt{TOF STICK} in
section~\ref{subsec:popup-profile}.
\end{itemize}

\subsubsection*{Data}
\begin{itemize}
\setlength{\itemsep}{0pt}
\setlength{\parsep}{0pt}
\item $H,V,\delta$ come from the \texttt{data-popup-field} controls
\texttt{releaseAltitude}, \texttt{releaseSpeed} and \texttt{diveAngle}
in \path{templates/dtc.html}. The \texttt{rangeModel} control selects
\texttt{low\_drag}, \texttt{high\_drag} or \texttt{manual}.
These are pilot inputs, not values extracted from a BMS weapon file.
\item $c_V$ is declared in \path{templates/dtc.html}.
$\Delta t,k,g,A_d,\epsilon$ are literals in
\path{templates/dtc_popup.js:computeDragBombTrajectory}.
The approximation has no wind or atmospheric-density input.
\item $B,t_f$ are returned by \texttt{computeDragBombTrajectory}.
\texttt{computePopupPlan} retains the unrounded range for geometry
and writes a whole-foot value into the range input.
Manual mode reads $B$ from \texttt{bombRange} and leaves $t_f$
unavailable. A failed trajectory or nonpositive range rejects the plan.
\end{itemize}

\subsubsection*{Example}
The input fixture in \path{tests/popup_briefing.test.cjs} uses
$H=2000$ ft, $V=520$ kt, $\delta=15^\circ$ and low drag:
$u_0=847.76$, $w_0=227.16$ ft/s.
The first step gives $a_0=-23.10000$ ft/s$^2$,
$x_1=84.66050$ ft and $h_1=1977.12312$ ft.
Ground crossing occurs at $m=62$:
\[
 h_{62}=-26.83394,\quad \Delta h_{61}=-42.504855,\quad
 x_{62}=4812.13003,\quad \Delta x_{61}=70.56950\ {\rm ft}.
\]
Thus $f=0.63131471$, $B=4767.57847$ ft and $t_f=6.13687$ s;
the displayed range is \texttt{4768}. The fixture's manual range
5138 ft is unused in this mode.

\subsubsection*{Code}
\path{templates/dtc_popup.js}, \texttt{computeDragBombTrajectory}.
The guard rejects a trajectory reaching 10000 iterations.
\texttt{altitude}/\texttt{lastStepAltitudeChange} gives $f$:
\begin{lstlisting}[numbers=left,firstnumber=299,numberstyle=\tiny,xleftmargin=2em]
function computeDragBombTrajectory(releaseAltitude, releaseSpeed, diveAngle, rangeModel) {
  const stepSeconds = 0.1;
  const dragScale = rangeModel === "high_drag" ? 1.0 : 0.165;
  const gravity = rangeModel === "high_drag" ? 20.91505 : 32.177;
  const diveRadians = diveAngle * Math.PI / 180;
  let horizontalSpeed = Math.round(Math.cos(diveRadians) * releaseSpeed * FEET_PER_KNOT_SECOND * 100) / 100;
  let verticalSpeedDown = Math.round(Math.sin(diveRadians) * releaseSpeed * FEET_PER_KNOT_SECOND * 100) / 100;
  let rangeFeet = 0;
  let altitude = releaseAltitude;
  let lastStepRange = 0;
  let lastStepAltitudeChange = 0;
  let timeSeconds = 0;
  let guard = 0;
  while (altitude > 0 && guard < 10000) {
    guard += 1;
    const horizontalDrag = dragScale * 140 * Math.abs(horizontalSpeed) / Math.sqrt(horizontalSpeed * horizontalSpeed + 0.1);
    const acceleration = horizontalSpeed > 0 ? -horizontalDrag : 0;
    lastStepRange = horizontalSpeed * stepSeconds + 0.5 * acceleration * stepSeconds * stepSeconds;
    lastStepAltitudeChange = -verticalSpeedDown * stepSeconds - 0.5 * gravity * stepSeconds * stepSeconds;
    rangeFeet += lastStepRange;
    altitude += lastStepAltitudeChange;
    timeSeconds += stepSeconds;
    horizontalSpeed = Math.max(0, horizontalSpeed + acceleration * stepSeconds);
    verticalSpeedDown += gravity * stepSeconds;
  }
  if (guard >= 10000 || !Number.isFinite(altitude) || !Number.isFinite(lastStepAltitudeChange) || lastStepAltitudeChange === 0) return null;
  const overshootFraction = altitude / lastStepAltitudeChange;
  return {
    rangeFeet: rangeFeet - overshootFraction * lastStepRange,
    timeSeconds: timeSeconds - overshootFraction * stepSeconds,
  };
}
\end{lstlisting}

\subsection{Stick, tracking and climb profile}\label{subsec:popup-profile}
\subsubsection*{Values}
The calculator and Pop-up briefing display \texttt{TOF STICK} (seconds
to one decimal), \texttt{CA} (climb angle), \texttt{PDWN}/\texttt{APEX}
(altitudes), and \texttt{AOD}, \texttt{PUP}, \texttt{RADIUS}, \texttt{MAP}
(distances in whole feet). Support values include \texttt{STICK LEN},
\texttt{STICK DUR}, \texttt{RELEASE RANGE}, \texttt{GROUND SPEED},
\texttt{TRACKING}, \texttt{TRACK ALT} and \texttt{ANGLE OFF}.
AGL heights are also shown in MSL using section~\ref{subsec:terrain}.

\subsubsection*{Formulas}
\begin{align*}
 V_h&=V\cos\delta,& L&=(n-1)s,& t_s&=\frac{L}{c_VV_h},\\
 B_r&=B+\begin{cases}L/2,&\text{MID},\\0,&\text{START},\end{cases}&
 t_{\rm stick}&=t_f+t_s,\\
 X_{\rm track}&=c_GV_ht,& Z_{\rm track}&=c_GVt\sin\delta,&
 H_{\rm track}&=H+Z_{\rm track},\\
 MAP&=B_r+X_{\rm track},&
 AOD&=\frac{H}{\tan\delta}-B_r,\\
 \chi&=\begin{cases}\delta+5^\circ,&\delta\leq15^\circ,\\
                      \delta+10^\circ,&\delta>15^\circ,\end{cases}&
 \Omega&=2\chi,& R&=\frac{(c_GV)^2}{g_TG},\\
 H_{\rm PD}&=H_{\rm track}-R(\cos\delta-\cos\chi),&
 H_{\rm apex}&=H_{\rm PD}+R(1-\cos\chi),\\
 d_{\rm pop}&=\frac{H_{\rm PD}-H_{\rm in}}{\tan\chi}.
\end{align*}

\subsubsection*{Variables}
\begin{itemize}
\setlength{\itemsep}{0pt}
\setlength{\parsep}{0pt}
\interlinepenalty=10000
\item $H,V,\delta$: release height (ft AGL), TAS (knots) and dive
angle, as in section~\ref{subsec:popup-bomb}.
$B,t_f$ are its bomb range (ft) and fall time (s).
\item $V_h$: horizontal component of release TAS, in knots;
the app labels it \texttt{GROUND SPEED}. Wind is not included.
\item $n=\max(1,\lfloor n_{\rm input}+0.5\rfloor)$: integer pulse count,
from the entered $n_{\rm input}$. $s$ is spacing in ft;
$L$ is total stick length in ft and $t_s$ is ripple duration in seconds.
\item $B_r$: target-to-first-release distance in ft. MID adds half
the stick length to centre the stick on the target; START uses $B$.
$t_{\rm stick}$ is time from first release to final impact, in seconds.
\item $t$: entered tracking time in seconds.
$X_{\rm track},Z_{\rm track}$: horizontal distance and height lost
during tracking, in ft; $H_{\rm track}$: tracking-start height in ft AGL.
\item $MAP$: target-to-maneuver-aim-point distance in ft;
$AOD$: aim-off distance beyond the target, in ft.
\item $\chi,\Omega$: climb angle and ingress/attack angle-off, in degrees.
$G$ is the selected dimensionless G value: 3.5 or 5.
$R$ is the turn/pull-down radius in ft.
\item $H_{\rm PD},H_{\rm apex},H_{\rm in}$: pull-down, apex and
entered ingress heights in ft AGL. $d_{\rm pop}$ is the horizontal
PUP-to-pull-down distance in ft, displayed as \texttt{PUP}.
\item $c_V=1.68781$ and $c_G=1.69$: application conversion factors
in ft/s per knot, used for ripple timing and tracking/turn geometry
respectively. $g_T=32.2$ ft/s$^2$ is the fixed acceleration constant
in the radius formula.
\end{itemize}

\subsubsection*{Data}
\begin{itemize}
\setlength{\itemsep}{0pt}
\setlength{\parsep}{0pt}
\item $H,V,\delta,H_{\rm in},t,n_{\rm input},s$ come from
\texttt{releaseAltitude}, \texttt{releaseSpeed}, \texttt{diveAngle},
\texttt{ingressAltitude}, \texttt{timeOnFinal}, \texttt{ripplePulses}
and \texttt{stickSpacing} in \path{templates/dtc.html}.
\texttt{stickMode} selects MID/START.
\item $G$ is selected by \texttt{gProfile}: the UI labels
\texttt{3-3.5 G} and \texttt{4.5-5 G} correspond to fixed values
3.5 and 5 in \path{templates/dtc_popup.js:computePopupPlan}.
The same function contains $c_G,g_T$ and the climb-angle rule;
$c_V$ is declared in \path{templates/dtc.html}.
\item $B,t_f$ come from section~\ref{subsec:popup-bomb};
manual range mode leaves $t_{\rm stick}$ unavailable, displayed
as \texttt{N/A}. The remaining values are computed in
\texttt{computePopupPlan} and formatted by \texttt{formatPopupOutput}
and \texttt{formatPopupBriefing} in \path{templates/dtc_popup.js}.
The app rejects $AOD\leq0$ or $d_{\rm pop}\leq0$.
\end{itemize}

\subsubsection*{Example}
Continuing the fixture in \path{tests/popup_briefing.test.cjs},
$n=2$, $s=100$ ft, MID, $t=5$ s, $H_{\rm in}=500$ ft and $G=3.5$:
\begin{align*}
 V_h&=520\cos15^\circ=502.28143\ {\rm kt},&
 L&=(2-1)100=100\ {\rm ft},\\
 t_s&=\frac{100}{1.68781(502.28143)}=0.11796\ {\rm s},&
 B_r&=4767.57847+50=4817.57847\ {\rm ft},\\
 t_{\rm stick}&=6.13687+0.11796=6.25483\ {\rm s},&
 MAP&=4817.57847+4244.27808=9061.85655\ {\rm ft}.
\end{align*}
The same calculation gives $AOD=2646.52314$ ft,
$\chi=20^\circ$, $\Omega=40^\circ$, $R=6852.61260$ ft,
$H_{\rm track}=3137.25088$ ft, $H_{\rm PD}=2957.48489$ ft,
$H_{\rm apex}=3370.74800$ ft and $d_{\rm pop}=6751.88424$ ft.
The briefing shows \texttt{TOF STICK 6.3 sec}, \texttt{CA 20 deg},
\texttt{PDWN 2957}, \texttt{APEX 3371}, \texttt{AOD 2647},
\texttt{PUP 6752}, \texttt{RADIUS 6853}, \texttt{MAP 9062} (feet).
Terrain is absent in this fixture, so MSL labels remain pending.

\subsubsection*{Code}
\path{templates/dtc_popup.js}, \texttt{computePopupPlan};
\texttt{ripplePulses} has already been rounded and clamped to at least 1:
\begin{lstlisting}[numbers=left,firstnumber=406,numberstyle=\tiny,xleftmargin=2em]
const diveRadians = diveAngle * Math.PI / 180;
const groundSpeed = releaseSpeed * Math.cos(diveRadians);
const groundSpeedFps = groundSpeed * FEET_PER_KNOT_SECOND;
const stickLength = Math.max(0, ripplePulses - 1) * stickSpacing;
const stickHalfLength = stickMode === "mid" ? stickLength / 2 : 0;
const stickDuration = groundSpeedFps > 0 ? stickLength / groundSpeedFps : NaN;
const releaseReferenceRange = bombRange + stickHalfLength;
const stickTimeOfFall = Number.isFinite(bombTimeOfFall) && Number.isFinite(stickDuration)
  ? bombTimeOfFall + stickDuration
  : NaN;
const horizontalTrackingDistance = groundSpeed * 1.69 * timeOnFinal;
const verticalTrackingDistance = releaseSpeed * 1.69 * timeOnFinal * Math.sin(diveRadians);
const mapDistance = releaseReferenceRange + horizontalTrackingDistance;
const actionRangeFeet = actionRangeNm * FEET_PER_NM;
const aod = releaseAltitude / Math.tan(diveRadians) - releaseReferenceRange;
if (aod <= 0) return { error: "AOD is negative for this profile. Increase release altitude or reduce bomb range." };
const climbAngle = diveAngle <= 15 ? diveAngle + 5 : diveAngle + 10;
const climbRadians = climbAngle * Math.PI / 180;
const angleOff = 2 * climbAngle;
const turnG = gProfile === "high" ? 5 : 3.5;
const trackAltitude = releaseAltitude + verticalTrackingDistance;
const turnRadius = ((releaseSpeed * 1.69) ** 2) / (32.2 * turnG);
\end{lstlisting}

Later in the same function, the tracking height and radius give the
pull-down/apex heights and PUP-to-pull-down distance:
\begin{lstlisting}[numbers=left,firstnumber=436,numberstyle=\tiny,xleftmargin=2em]
const pullDownAltitude = trackAltitude - turnRadius * (Math.cos(diveRadians) - Math.cos(climbRadians));
const apexAltitude = pullDownAltitude + turnRadius * (1 - Math.cos(climbRadians));
const popToPullDownDistance = (pullDownAltitude - ingressAltitude) / Math.tan(climbRadians);
if (popToPullDownDistance <= 0) return { error: "Ingress altitude is above the computed pull-down altitude." };
\end{lstlisting}

\subsection{Recovery height and clearance}\label{subsec:popup-recovery}
\subsubsection*{Values}
\texttt{ALTLOSS}, \texttt{GNDC} and \texttt{MRA}, in whole feet.
\texttt{GNDC} is AGL; \texttt{MRA} is shown in AGL and MSL.
Support values show \texttt{DELAY LOSS}, \texttt{PULL LOSS},
\texttt{PULL RADIUS}, \texttt{MRA BUFFER}, and the recovery G/delay.

\subsubsection*{Formulas}
\begin{align*}
 t_d&=t_s+\tau_d,& R_r&=\frac{(c_VV)^2}{g_TG_r},\\
 \Delta H_d&=c_VV\sin\delta\,t_d,&
 \Delta H_r&=R_r(1-\cos\delta),\\
 \Delta H&=\Delta H_d+\Delta H_r,&
 H_{\rm clear}&=H-\Delta H,&
 H_{\rm MRA}&=\Delta H+H_b.
\end{align*}

\subsubsection*{Variables}
\begin{itemize}
\setlength{\itemsep}{0pt}
\setlength{\parsep}{0pt}
\interlinepenalty=10000
\item $H,V,\delta,t_s$: release height (ft AGL), release TAS (knots),
dive angle and ripple duration (s), from section~\ref{subsec:popup-profile}.
\item $\tau_d=2$ s: fixed delay after the last release; $t_d$:
total delay from the first release to the recovery pull.
\item $G_r=5$: fixed recovery G value, independent of the selected
turn profile. $R_r$ is the recovery pull radius in ft.
$c_V=1.68781$ ft/s per knot and $g_T=32.2$ ft/s$^2$ retain the
application constants above.
\item $\Delta H_d,\Delta H_r,\Delta H$: height lost during the delay,
during the pull to level, and their sum, in ft.
\item $H_{\rm clear}$: remaining height above ground after recovery,
in ft; $H_{\rm MRA}$: release height needed by this model to leave
the buffer $H_b=120$ ft. MSL conversion follows
section~\ref{subsec:terrain}.
\end{itemize}

\subsubsection*{Data}
\begin{itemize}
\setlength{\itemsep}{0pt}
\setlength{\parsep}{0pt}
\item $H,V,\delta$ come from the same DTC pop-up input controls
as above; $t_s$ is computed from pulse count and spacing in
\path{templates/dtc_popup.js:computePopupPlan}.
\item $\tau_d,G_r,g_T$ are fixed literals in \texttt{computePopupPlan}.
$H_b$ is \texttt{POPUP\_MRA\_BUFFER\_FEET}, and $c_V$ is
\texttt{FEET\_PER\_KNOT\_SECOND}, both declared in
\path{templates/dtc.html}. These parameters are not selected
from aircraft-specific recovery tables.
\item $\Delta H,H_{\rm clear},H_{\rm MRA}$ are computed by
\texttt{computePopupPlan}; \texttt{formatPopupOutput} and
\texttt{formatPopupBriefing} display them.
The arithmetic assumes a common ground reference through the recovery;
it does not sample terrain along a recovery path.
\end{itemize}

\subsubsection*{Example}
For the same \path{tests/popup_briefing.test.cjs} fixture:
\begin{align*}
 t_d&=0.11795852+2=2.11795852\ {\rm s},\\
 R_r&=\frac{(1.68781\cdot520)^2}{32.2\cdot5}=4784.40486\ {\rm ft},\\
 \Delta H_d&=1.68781\cdot520\sin15^\circ\cdot2.11795852
              =481.10579\ {\rm ft},\\
 \Delta H_r&=4784.40486(1-\cos15^\circ)=163.02464\ {\rm ft}.
\end{align*}
Hence $\Delta H=644.13043$ ft,
$H_{\rm clear}=2000-644.13043=1355.86957$ ft and
$H_{\rm MRA}=644.13043+120=764.13043$ ft.
The briefing displays \texttt{ALTLOSS 644 ft},
\texttt{GNDC 1356 ft AGL}, and \texttt{MRA 764 ft AGL}
(with MSL pending in this fixture).

\subsubsection*{Code}
\path{templates/dtc_popup.js}, \texttt{computePopupPlan}, delay and
pull contributions:
\begin{lstlisting}[numbers=left,firstnumber=428,numberstyle=\tiny,xleftmargin=2em]
const recoveryPullDelay = stickDuration + 2;
const recoveryPullG = 5;
const releaseSpeedFps = releaseSpeed * FEET_PER_KNOT_SECOND;
const recoveryDelayAltitudeLoss = releaseSpeedFps * Math.sin(diveRadians) * recoveryPullDelay;
const recoveryPullRadius = (releaseSpeedFps ** 2) / (32.2 * recoveryPullG);
const recoveryPullAltitudeLoss = recoveryPullRadius * (1 - Math.cos(diveRadians));
const recoveryAltitudeLoss = recoveryDelayAltitudeLoss + recoveryPullAltitudeLoss;
const recoveryGroundClearance = releaseAltitude - recoveryAltitudeLoss;
\end{lstlisting}

The same function adds the buffer when returning the plan values:
\begin{lstlisting}[numbers=left,firstnumber=533,numberstyle=\tiny,xleftmargin=2em]
recoveryMraAltitude: recoveryAltitudeLoss + POPUP_MRA_BUFFER_FEET,
recoveryMraBuffer: POPUP_MRA_BUFFER_FEET,
\end{lstlisting}

\subsection{Map points, turn and navigation offsets}\label{subsec:popup-map}
\subsubsection*{Values}
The map shows action, PUP, pull-down, MAP, target and OA1 points
and the connecting turn. \texttt{POP BRG}/\texttt{POP RNG} give the
PUP-to-target bearing and range (degrees to one decimal, NM to one
decimal). Support values show point-to-point distances, headings,
and \texttt{PUP OFFSET}/\texttt{OA1 OFFSET} bearing/range pairs
(true degrees to one decimal and whole feet). The action-point tooltip
also shows the \texttt{Offset} angle between ingress and the leg to PDP.

\subsubsection*{Formulas}
The entered heading determines either $\theta_A$ or $\theta_I$;
the other follows from the selected side:
\[
 \theta_I=(\theta_A-\eta\Omega)\bmod360^\circ,\qquad
 \theta_A=(\theta_I+\eta\Omega)\bmod360^\circ.
\]
Construct the action point, MAP, release point and aim-off point:
\begin{align*}
 \mathbf P&=\mathbf T-d_A\mathbf e(\theta_I),&
 \mathbf M&=\mathbf T-MAP\,\mathbf e(\theta_A),\\
 \mathbf X&=\mathbf T-B_r\mathbf e(\theta_A),&
 \mathbf O&=\mathbf T+AOD\,\mathbf e(\theta_A).
\end{align*}
For the circle tangent to the attack heading at MAP:
\[
 \mathbf C=\mathbf M+\eta R\mathbf n(\theta_A),\qquad
 \mathbf v=\mathbf P-\mathbf C,\qquad
 \mathbf Q_\pm=\mathbf C+\frac{R^2}{\|\mathbf v\|^2}\mathbf v
       \ \pm\frac{R\sqrt{\|\mathbf v\|^2-R^2}}{\|\mathbf v\|^2}J\mathbf v.
\]
The selected pull-down point $\mathbf Q$ is the candidate whose
computed arc endpoint is closest to $\mathbf M$. Then
\[
 \mathbf U=\mathbf Q-d_{\rm pop}\mathbf e(\theta_o),\qquad
 D_{\rm POP}=\frac{\|\mathbf T-\mathbf U\|}{c_N},\qquad
 \omega=\min(\alpha,360^\circ-\alpha).
\]
Bearings and the remaining straight-line distances use the formulas
in section~\ref{subsec:offset-geometry}; magnetic display uses
section~\ref{subsec:magnetic-headings}.

\subsubsection*{Variables}
\begin{itemize}
\setlength{\itemsep}{0pt}
\setlength{\parsep}{0pt}
\interlinepenalty=10000
\item All bold point vectors are $(N,E)$ in feet.
$\mathbf S,\mathbf T$: selected reference steerpoint and target.
$\mathbf P,\mathbf M,\mathbf X,\mathbf O,\mathbf Q,\mathbf U$:
action, MAP, first release, OA1, pull-down (PDP), and PUP positions.
\item $\theta_A,\theta_I,\theta_o$: true attack, ingress and
action-to-PDP headings, in degrees.
$\eta=+1$ for a right turn and $-1$ for a left turn.
$\Omega$ is the angle-off from section~\ref{subsec:popup-profile}.
\item $\mathbf e(\theta)=(\cos\theta,\sin\theta)$: forward unit vector;
$\mathbf n(\theta)=(-\sin\theta,\cos\theta)$: right unit normal.
The arguments are converted to radians for evaluation.
\item $d_A=c_Nr_A$: action-to-target distance in ft, from entered
range $r_A$ in NM; $c_N=6076.12$ ft/NM is \texttt{FEET\_PER\_NM}.
$MAP,B_r,AOD,R,d_{\rm pop}$ are the distances/radius from
section~\ref{subsec:popup-profile}, in ft.
\item $\mathbf C$: turn-circle centre at MAP; $\mathbf v$: vector
from that centre to the action point.
$J(v_N,v_E)=(v_E,-v_N)$ is a quarter-turn operator.
$\|\cdot\|$ is Euclidean length. The two tangent candidates
$\mathbf Q_+,\mathbf Q_-$ require $\|\mathbf v\|>R$.
\item $j\in\{+,-\}$: tangent-candidate index. For each candidate,
$\theta_j$ is the true bearing from $\mathbf P$ to $\mathbf Q_j$,
rounded to $0.1^\circ$ by \texttt{bearingRangeBetween}.
The arc is reconstructed from that rounded heading:
\begin{align*}
 \mathbf C_j&=\mathbf Q_j+\eta R\mathbf n(\theta_j),\\
 \psi_j&=\eta\bigl[(\eta(\theta_A-\theta_j))\bmod360^\circ\bigr],\\
 \mathbf A_j&=\mathbf C_j+\mathcal R_{\psi_j}(\mathbf Q_j-\mathbf C_j).
\end{align*}
$\mathbf C_j$ is the reconstructed centre, $\psi_j$ the signed
turn angle, and $\mathbf A_j$ the arc endpoint.
The rotation is
$\mathcal R_\psi(N,E)=(N\cos\psi-E\sin\psi,\,
N\sin\psi+E\cos\psi)$.
For $|\psi_j|<0.1^\circ$, the code returns $\mathbf A_j=\mathbf Q_j$.
\item $j_*=\arg\min_{j\in\{+,-\}}\|\mathbf A_j-\mathbf M\|$:
chosen candidate; $\mathbf Q=\mathbf Q_{j_*}$ and
$\theta_o=\theta_{j_*}$. This is an endpoint-error comparison;
the code applies no maximum-error threshold.
\item $\alpha=(\theta_o-\theta_I)\bmod360^\circ$:
wrapped difference between the offset-leg and ingress headings.
$\omega$ is the smaller angle between them, displayed as \texttt{Offset}.
\item $D_{\rm POP}$: PUP-to-target range in NM.
The PUP and OA1 offset rows are measured from $\mathbf S$ to
$\mathbf U$ and $\mathbf O$, respectively.
\end{itemize}

\subsubsection*{Data}
\begin{itemize}
\setlength{\itemsep}{0pt}
\setlength{\parsep}{0pt}
\item $\mathbf S$ is read from the selected
\texttt{[STPT].target\_k} in the loaded \path{<callsign/save>.ini},
through \texttt{getBmsCoordForStpt} in \path{templates/dtc.html}.
VRP mode uses $\mathbf T=\mathbf S$.
VIP mode projects $\mathbf T$ from $\mathbf S$ using the
\texttt{VIP} bearing/range in \texttt{[NAV OFFSETS]} or its edited
DTC fields, as in section~\ref{subsec:offset-geometry}.
\item $r_A,\eta$ and the entered heading come from
\texttt{actionRangeNm}, \texttt{side} and \texttt{heading};
\texttt{headingMode} selects attack or ingress. These are controls
in \path{templates/dtc.html}. The entered heading is converted
to true using the selected point's variation as described in
section~\ref{subsec:magnetic-headings}. $c_N$ is declared in
\path{templates/dtc.html}.
\item $\mathbf C,\mathbf Q_\pm,\theta_j,\mathbf A_j,j_*$ are computed by
\texttt{getTurnCenterAtFinalPoint}, \texttt{getCircleTangentsFromPoint},
\texttt{bearingRangeBetween}, \texttt{buildTurnArcCoords} and
\texttt{solvePopupTurnGeometry} in \path{templates/dtc.html}.
The chosen tangent need not land exactly at MAP after heading rounding.
\item $\mathbf P,\mathbf M,\mathbf X,\mathbf O,\mathbf U$ and their
bearing/range pairs are computed by
\path{templates/dtc_popup.js:computePopupPlan}. Missing target data,
no tangent, or $d_{\rm pop}>\|\mathbf Q-\mathbf P\|$ rejects the plan.
\texttt{applyPopupPlan} writes the computed offsets into the selected
\texttt{OA1} row and \texttt{VIPPUP}/\texttt{VRPPUP}.
\item $D_{\rm POP}$ and the displayed point distances use unrounded
coordinates; bearings first pass through the $0.1^\circ$ DTC
formatter. The support offset rows retain true bearings, while
the pop-up heading labels apply the magnetic conversion.
\end{itemize}

\subsubsection*{Example}
The same \path{tests/popup_briefing.test.cjs} fixture supplies the
synthetic reference $\mathbf S=(100000,100000)$ ft, VRP mode,
$r_A=4.5$ NM, a right turn and attack heading $130^\circ$ magnetic.
Its mocked variation is $-8^\circ$, giving
$\theta_A=122^\circ$, $\theta_I=82^\circ$ and $d_A=27342.54$ ft.
Together with the profile above:
\[
 \mathbf P=(96194.65,72923.56),\qquad
 \mathbf M=(104802.05,92315.11),\qquad
 \mathbf C=(98990.71,88683.78)\ {\rm ft}.
\]
The candidate arc endpoint errors are 13705.23 ft and 1.02 ft.
The second candidate is selected:
\[
 \mathbf Q=(104575.87,84713.36)\ {\rm ft},\qquad
 \theta_o=54.6^\circ,\qquad
 \mathbf U=(100664.63,79209.71)\ {\rm ft}.
\]
The briefing shows \texttt{POP BRG 99.8 deg M} and
\texttt{POP RNG 3.4 NM}; the action-point \texttt{Offset} is
$\omega=|54.6-82|=27.4^\circ$. Support offsets are
\texttt{PUP BRG 271.8 / RNG 20801 ft} and
\texttt{OA1 BRG 122 / RNG 2647 ft}, both true.
These coordinates and variation are test inputs, not BMS-file samples.

\subsubsection*{Code}
\path{templates/dtc.html}, \texttt{getTurnCenterAtFinalPoint},
the right-normal construction of $\mathbf C$:
\begin{lstlisting}[numbers=left,firstnumber=816,numberstyle=\tiny,xleftmargin=2em]
const radians = normalizeBearing(finalHeading) * Math.PI / 180;
const rightNormal = {
  northFeet: -Math.sin(radians),
  eastFeet: Math.cos(radians),
};
return turnDirection === "left"
  ? {
    northFeet: finalCoord.northFeet - rightNormal.northFeet * radiusFeet,
    eastFeet: finalCoord.eastFeet - rightNormal.eastFeet * radiusFeet,
  }
  : {
    northFeet: finalCoord.northFeet + rightNormal.northFeet * radiusFeet,
    eastFeet: finalCoord.eastFeet + rightNormal.eastFeet * radiusFeet,
  };
\end{lstlisting}

The same file's \texttt{getCircleTangentsFromPoint} produces
$\mathbf Q_\pm$; \texttt{solvePopupTurnGeometry} ranks the resulting
arc endpoints. Here \texttt{dx} is east and \texttt{dy} is north:
\begin{lstlisting}[numbers=left,firstnumber=832,numberstyle=\tiny,xleftmargin=2em]
function getCircleTangentsFromPoint(pointCoord, centerCoord, radiusFeet) {
  if (!pointCoord || !centerCoord || !Number.isFinite(radiusFeet) || radiusFeet <= 0) return [];
  const dx = pointCoord.eastFeet - centerCoord.eastFeet;
  const dy = pointCoord.northFeet - centerCoord.northFeet;
  const distanceSquared = dx * dx + dy * dy;
  if (distanceSquared <= radiusFeet * radiusFeet) return [];
  const root = Math.sqrt(distanceSquared - radiusFeet * radiusFeet);
  return [1, -1].map((sign) => {
    const eastFeet = centerCoord.eastFeet + (radiusFeet * radiusFeet * dx + sign * radiusFeet * root * -dy) / distanceSquared;
    const northFeet = centerCoord.northFeet + (radiusFeet * radiusFeet * dy + sign * radiusFeet * root * dx) / distanceSquared;
    return { northFeet, eastFeet };
  });
}

function solvePopupTurnGeometry(actionCoord, mapCoord, attackHeading, turnRadius, turnDirection) {
  const centerCoord = getTurnCenterAtFinalPoint(mapCoord, attackHeading, turnRadius, turnDirection);
  if (!centerCoord) return null;
  const candidates = getCircleTangentsFromPoint(actionCoord, centerCoord, turnRadius)
    .map((pdpCoord) => {
      const offsetLegMetrics = bearingRangeBetween(actionCoord, pdpCoord);
      if (!offsetLegMetrics) return null;
      const offsetLegHeading = Number.parseFloat(offsetLegMetrics.brg);
      const turnArcCoords = buildTurnArcCoords(pdpCoord, offsetLegHeading, attackHeading, turnRadius, turnDirection);
      const arcEndCoord = turnArcCoords.length ? turnArcCoords[turnArcCoords.length - 1] : null;
      return {
        pdpCoord,
        offsetLegMetrics,
        offsetLegHeading,
        turnArcCoords,
        errorFeet: coordDistanceFeet(arcEndCoord, mapCoord),
      };
    })
    .filter((candidate) => candidate && Number.isFinite(candidate.errorFeet));
  if (!candidates.length) return null;
  candidates.sort((a, b) => a.errorFeet - b.errorFeet);
  return candidates[0];
}
\end{lstlisting}

\texttt{buildTurnArcCoords} in the same file reconstructs the circle
from the rounded starting heading and rotates the initial radius.
At \texttt{index === steps}, the appended point is $\mathbf A_j$;
\texttt{getTurnDelta} computes $\psi_j$:
\begin{lstlisting}[numbers=left,firstnumber=778,numberstyle=\tiny,xleftmargin=2em]
const turnDelta = getTurnDelta(startHeading, endHeading, turnDirection);
if (!Number.isFinite(turnDelta) || Math.abs(turnDelta) < 0.1) return [startCoord];
const startRadians = normalizeBearing(startHeading) * Math.PI / 180;
const rightNormal = {
  northFeet: -Math.sin(startRadians),
  eastFeet: Math.cos(startRadians),
};
const centerCoord = turnDelta >= 0
  ? {
    northFeet: startCoord.northFeet + rightNormal.northFeet * radiusFeet,
    eastFeet: startCoord.eastFeet + rightNormal.eastFeet * radiusFeet,
  }
  : {
    northFeet: startCoord.northFeet - rightNormal.northFeet * radiusFeet,
    eastFeet: startCoord.eastFeet - rightNormal.eastFeet * radiusFeet,
  };
const relative = {
  northFeet: startCoord.northFeet - centerCoord.northFeet,
  eastFeet: startCoord.eastFeet - centerCoord.eastFeet,
};
const steps = Math.max(8, Math.ceil(Math.abs(turnDelta) / 6));
const points = [];
for (let index = 0; index <= steps; index += 1) {
  const radians = (turnDelta * index / steps) * Math.PI / 180;
  const cos = Math.cos(radians);
  const sin = Math.sin(radians);
  const eastFeet = relative.eastFeet * cos + relative.northFeet * sin;
  const northFeet = -relative.eastFeet * sin + relative.northFeet * cos;
  points.push({
    northFeet: centerCoord.northFeet + northFeet,
    eastFeet: centerCoord.eastFeet + eastFeet,
  });
\end{lstlisting}

\path{templates/dtc_popup.js}, \texttt{computePopupPlan}, converts
the entered heading, constructs the points, and measures the offsets.
\texttt{offsetLegHeading} is the selected rounded $\theta_o$:
\begin{lstlisting}[numbers=left,firstnumber=441,numberstyle=\tiny,xleftmargin=2em]
const magneticVariation = getPopupMagneticVariation();
// Convert the pilot's heading once; geometry and DTC offset rows stay true.
const trueHeading = normalizeBearing(heading + (magneticVariation ?? 0));
let attackHeading;
let ingressHeading;
if (headingMode === "attack") {
  attackHeading = trueHeading;
  ingressHeading = normalizeBearing(attackHeading + (side === "right" ? -angleOff : angleOff));
} else {
  ingressHeading = trueHeading;
  attackHeading = normalizeBearing(ingressHeading + (side === "right" ? angleOff : -angleOff));
}

const turnDirection = side === "left" ? "left" : "right";
const actionCoord = projectBearingRange(targetCoord, ingressHeading + 180, actionRangeFeet);
const mapCoord = projectBearingRange(targetCoord, attackHeading + 180, mapDistance);
const releaseCoord = projectBearingRange(targetCoord, attackHeading + 180, releaseReferenceRange);
const turnGeometry = solvePopupTurnGeometry(actionCoord, mapCoord, attackHeading, turnRadius, turnDirection);
if (!turnGeometry) return { error: "Could not solve action-to-MAP turn geometry." };
const pdpCoord = turnGeometry.pdpCoord;
const offsetLegMetrics = turnGeometry.offsetLegMetrics;
const offsetLegHeading = turnGeometry.offsetLegHeading;
const offsetLegDistance = coordDistanceFeet(actionCoord, pdpCoord);
if (popToPullDownDistance > offsetLegDistance) {
  return { error: "PUP would fall before the action point. Increase action range or adjust the vertical profile." };
}
const pupCoord = projectBearingRange(pdpCoord, offsetLegHeading + 180, popToPullDownDistance);
const oaCoord = projectBearingRange(targetCoord, attackHeading, aod);
if (!actionCoord || !pdpCoord || !pupCoord || !oaCoord) return { error: "Could not compute pop-up geometry." };
const turnArcCoords = turnGeometry.turnArcCoords;
const rolloutCoord = mapCoord;

const pupMetrics = bearingRangeBetween(selectedCoord, pupCoord);
const oaMetrics = bearingRangeBetween(selectedCoord, oaCoord);
if (!pupMetrics || !oaMetrics) return { error: "Could not compute NAV offset rows." };
const ingressToOffset = Math.abs(normalizeBearing(offsetLegHeading - ingressHeading));
const offsetAngle = ingressToOffset > 180 ? 360 - ingressToOffset : ingressToOffset;
\end{lstlisting}

\begin{thebibliography}{5}
\bibitem{meeus1991}
Jean Meeus, \emph{Astronomical Algorithms}, Willmann--Bell, 1991.
\bibitem{ussa1976}
NOAA, NASA and U.S. Air Force, \emph{U.S. Standard Atmosphere, 1976},
NASA-TM-X-74335, 1976, Part~1.
\url{https://ntrs.nasa.gov/citations/19770009539}.
\bibitem{aiken1946}
William S. Aiken, Jr., \emph{Standard Nomenclature for Airspeeds with
Tables and Charts for Use in Calculation of Airspeed},
NACA Report~837, 1946.
\url{https://ntrs.nasa.gov/api/citations/19930091914/downloads/19930091914.pdf}.
\bibitem{how2004}
Jonathan P. How, \emph{Aircraft Stability and Control},
MIT OpenCourseWare 16.333, Fall 2004, Lecture~1:
Equilibrium States, Aircraft Performance, Introduction to Basic Terms.
\url{https://ocw.mit.edu/courses/16-333-aircraft-stability-and-control-fall-2004/9befd7ea465e31f9ff763c5587815264_lecture_1.pdf}.
\bibitem{oakdesign_converter}
oakdesign, \emph{Falcon BMS Tacview Converter}, software repository,
\path{src/utils/coordinate_converter.py}
(\texttt{create\_proj\_string} and \texttt{CoordinateConverter}).
\url{https://github.com/oakdesign/falcon-bms-tacview-converter/}.
\end{thebibliography}
\end{document}
